% ============================================================
%  Smart Farm AI v3.0 — Rapport de Fin de Cycle d'Ingénierie
%  Tek-Up 2024-2025 | Mohamed Sayari
%  FICHIER UNIQUE — compiler avec : pdflatex rapport_complet.tex
% ============================================================
\documentclass[12pt,a4paper,oneside]{report}

%% ── Encodage & langue ───────────────────────────────────────
\usepackage[utf8]{inputenc}
\usepackage[T1]{fontenc}
\usepackage[french]{babel}

%% ── Mise en page ────────────────────────────────────────────
\usepackage[left=3cm,right=2.5cm,top=2.5cm,bottom=2.5cm,headheight=15pt]{geometry}
\usepackage{setspace}\onehalfspacing
\usepackage{lmodern}
\usepackage{microtype}

%% ── Maths ───────────────────────────────────────────────────
\usepackage{amsmath,amssymb}

%% ── Graphiques ──────────────────────────────────────────────
\usepackage{graphicx,xcolor,float,subcaption}

%% ── Couleurs ────────────────────────────────────────────────
\definecolor{farmgreen}{RGB}{34,139,34}
\definecolor{farmblue}{RGB}{0,71,171}
\definecolor{farmgray}{RGB}{100,100,100}
\definecolor{farmorange}{RGB}{230,126,34}
\definecolor{farmred}{RGB}{192,57,43}
\definecolor{codebg}{RGB}{248,248,248}
\definecolor{codeframe}{RGB}{200,200,200}

%% ── TikZ ────────────────────────────────────────────────────
\usepackage{tikz}
\usepackage{pgfplots}\pgfplotsset{compat=1.18}
\usetikzlibrary{shapes.geometric,shapes.multipart,arrows.meta,
  positioning,fit,backgrounds,calc,matrix,shadows.blur,
  decorations.pathreplacing,decorations.markings,patterns}

%% ── Tableaux ────────────────────────────────────────────────
\usepackage{booktabs,tabularx,longtable,multirow,colortbl,array}
\usepackage{pifont}

%% ── Listings ────────────────────────────────────────────────
\usepackage{listings}
\lstdefinestyle{pythonstyle}{
  language=Python,basicstyle=\ttfamily\footnotesize,
  backgroundcolor=\color{codebg},frame=single,rulecolor=\color{codeframe},
  keywordstyle=\color{farmblue}\bfseries,commentstyle=\color{farmgreen}\itshape,
  stringstyle=\color{farmred},numbers=left,numberstyle=\tiny\color{farmgray},
  stepnumber=1,numbersep=8pt,showstringspaces=false,breaklines=true,tabsize=4}
\lstdefinestyle{arducstyle}{
  language=C++,basicstyle=\ttfamily\footnotesize,
  backgroundcolor=\color{codebg},frame=single,rulecolor=\color{codeframe},
  keywordstyle=\color{farmblue}\bfseries,commentstyle=\color{farmgreen}\itshape,
  stringstyle=\color{farmred},numbers=left,numberstyle=\tiny\color{farmgray},
  stepnumber=1,numbersep=8pt,showstringspaces=false,breaklines=true}
\lstdefinestyle{jsonstyle}{
  basicstyle=\ttfamily\footnotesize,backgroundcolor=\color{codebg},
  frame=single,rulecolor=\color{codeframe},showstringspaces=false,
  breaklines=true,morestring=[b]",stringstyle=\color{farmred}}

%% ── Boîtes colorées ─────────────────────────────────────────
\usepackage[most]{tcolorbox}
\tcbuselibrary{skins,breakable}
\tcbset{
  technote/.style={enhanced,breakable,colback=farmblue!5,colframe=farmblue!60!black,
    fonttitle=\bfseries,coltitle=white,
    attach boxed title to top left={yshift=-2mm,xshift=4mm},
    boxed title style={colback=farmblue}},
  infobox/.style={enhanced,breakable,colback=farmgreen!8,colframe=farmgreen!60!black,
    fonttitle=\bfseries,coltitle=white,
    attach boxed title to top left={yshift=-2mm,xshift=4mm},
    boxed title style={colback=farmgreen}},
  warningbox/.style={enhanced,breakable,colback=farmorange!10,colframe=farmorange!70!black,
    fonttitle=\bfseries,coltitle=white,
    attach boxed title to top left={yshift=-2mm,xshift=4mm},
    boxed title style={colback=farmorange}}
}

%% ── En-têtes ────────────────────────────────────────────────
\usepackage{fancyhdr}
\pagestyle{fancy}\fancyhf{}
\fancyhead[L]{\small\itshape\leftmark}
\fancyhead[R]{\small\color{farmblue}\textbf{Smart Farm AI v3.0}}
\fancyfoot[C]{\small\thepage}
\fancyfoot[R]{\tiny\color{farmgray}Tek-Up 2024-2025}
\renewcommand{\headrulewidth}{0.4pt}
\renewcommand{\footrulewidth}{0.2pt}

%% ── Titres colorés ──────────────────────────────────────────
\usepackage{titlesec}
\titleformat{\chapter}[display]
  {\normalfont\huge\bfseries\color{farmblue}}
  {\chaptertitlename\ \thechapter}{18pt}{\Huge\color{farmgreen}}
\titleformat{\section}{\normalfont\Large\bfseries\color{farmblue}}{\thesection}{1em}{}
\titleformat{\subsection}{\normalfont\large\bfseries\color{farmgray}}{\thesubsection}{1em}{}

%% ── Liens PDF ───────────────────────────────────────────────
\usepackage[hidelinks,pdftitle={Smart Farm AI v3.0},pdfauthor={Mohamed Sayari}]{hyperref}

%% ── Espacement ──────────────────────────────────────────────
\setlength{\parskip}{6pt}\setlength{\parindent}{1.5em}

%% ── Macros ──────────────────────────────────────────────────
\newcommand{\farmAI}{\textbf{Smart Farm AI}\textsuperscript{TM}}
\newcommand{\api}[1]{\texttt{\color{farmblue}\small#1}}
\newcommand{\code}[1]{\texttt{\small\color{farmgray}#1}}
\newcommand{\yes}{\ding{52}}
\newcommand{\no}{\ding{56}}

%% ── Stickman TikZ (acteurs UML) ─────────────────────────────
\newcommand{\stickactor}[3]{%  x, y, label
  \begin{scope}[shift={(#1,#2)}]
    \draw[thick,farmgray](0,0.55)circle(0.18);
    \draw[thick,farmgray](0,0.37)--(0,-0.15);
    \draw[thick,farmgray](-0.22,0.15)--(0.22,0.15);
    \draw[thick,farmgray](0,-0.15)--(-0.18,-0.55);
    \draw[thick,farmgray](0,-0.15)--(0.18,-0.55);
  \end{scope}
  \node[font=\scriptsize\bfseries,align=center,text width=1.8cm]
    at (#1,#2-1.0){#3};
}

% ============================================================
\begin{document}
% ============================================================

%% ──────────────────────────────────────────────────────────────
%%  PAGE DE GARDE
%% ──────────────────────────────────────────────────────────────
\begin{titlepage}
\thispagestyle{empty}
\begin{center}
{\footnotesize\textbf{République Tunisienne}}\\[2pt]
{\footnotesize Ministère de l'Enseignement Supérieur et de la Recherche Scientifique}\\[12pt]

\begin{minipage}[c]{0.44\textwidth}\centering
\begin{tcolorbox}[enhanced,colback=farmblue!10,colframe=farmblue,arc=6pt,
  boxrule=1pt,width=4cm,height=2.2cm,valign=center,halign=center]
{\small\bfseries\color{farmblue}Tek-Up\\École d'Ingénieurs\\Tunis}
\end{tcolorbox}
\end{minipage}%
\begin{minipage}[c]{0.44\textwidth}\centering
\begin{tcolorbox}[enhanced,colback=farmgreen!10,colframe=farmgreen,arc=6pt,
  boxrule=1pt,width=4cm,height=2.2cm,valign=center,halign=center]
{\small\bfseries\color{farmgreen}Rapport de Fin\\de Cycle\\2024--2025}
\end{tcolorbox}
\end{minipage}

\vspace{0.5cm}
\rule{\textwidth}{1.5pt}\\[4pt]
{\large\textbf{École Privée d'Ingénieurs — Tek-Up}}\\
{\normalsize Spécialité : \textit{Data Science et Intelligence Artificielle}}\\[4pt]
\rule{\textwidth}{1.5pt}

\vspace{0.6cm}
{\large\textsc{Rapport de Fin de Cycle d'Ingénierie}}\\[4pt]
{\normalsize En vue de l'obtention du diplôme d'\\[2pt]
\textbf{Ingénieur en Data Science et Intelligence Artificielle}}

\vspace{0.7cm}
\begin{tcolorbox}[enhanced,colback=farmblue!8,colframe=farmblue,arc=8pt,
  boxrule=1.5pt,width=\textwidth]
\centering
{\LARGE\bfseries\color{farmblue}Smart Farm AI v3.0}\\[6pt]
{\large\color{farmgreen}Plateforme Agro-Technologique Multi-Tenant basée sur\\
l'Intelligence Artificielle, l'IoT et les LLM Souverains}\\[4pt]
{\normalsize\color{farmgray}\itshape
Détection automatisée des maladies agricoles, gestion apicole intelligente\\
et assistant conversationnel en Darija tunisien}
\end{tcolorbox}

\vspace{0.7cm}
\begin{minipage}{0.5\textwidth}\raggedright
{\normalsize\textbf{Réalisé par :}}\\[3pt]
{\large\color{farmblue}\textbf{Mohamed Sayari}}\\
{\small Classe : DSI-5}
\end{minipage}%
\begin{minipage}{0.5\textwidth}\raggedleft
{\normalsize\textbf{Encadré par :}}\\[3pt]
{\large\textbf{[Encadrant]}}\\
{\small Maître de conférences, Tek-Up}
\end{minipage}

\vfill
{\normalsize Année universitaire : \textbf{2024--2025}}
\end{center}
\end{titlepage}

%% ──────────────────────────────────────────────────────────────
%%  PAGES LIMINAIRES (numérotation romaine)
%% ──────────────────────────────────────────────────────────────
\pagenumbering{roman}

%% Dédicaces
\chapter*{Dédicaces}\thispagestyle{empty}
\addcontentsline{toc}{chapter}{Dédicaces}
\begin{center}\vspace{2cm}
{\large\itshape
À mes parents, pour leur soutien indéfectible.\\[1.2em]
À mes professeurs de Tek-Up,\\
qui m'ont transmis la passion pour la Data Science et l'IA.\\[1.2em]
À tous les agriculteurs tunisiens\\
dont le travail quotidien a inspiré ce projet.
}
\vspace{1.5cm}
\begin{flushright}\textit{Mohamed Sayari}\end{flushright}
\end{center}

\clearpage

%% Remerciements
\chapter*{Remerciements}\thispagestyle{empty}
\addcontentsline{toc}{chapter}{Remerciements}

Je tiens à exprimer ma profonde gratitude à toutes les personnes qui ont contribué à
la réalisation de ce projet de fin de cycle.

Mes sincères remerciements vont à mon encadrant pédagogique pour ses conseils avisés
et sa disponibilité tout au long de ce travail.

Je remercie également le corps enseignant de l'École Tek-Up pour la formation de
qualité dispensée en Data Science et Intelligence Artificielle, ainsi que l'UTAP et
l'AVFA pour leurs ressources documentaires sur l'agriculture tunisienne.

\clearpage

%% Résumé FR
\chapter*{Résumé}\thispagestyle{empty}
\addcontentsline{toc}{chapter}{Résumé}

\noindent\textbf{Titre :} Smart Farm AI v3.0 — Plateforme Agro-Technologique Multi-Tenant

\medskip
Ce rapport présente \farmAI{} v3.0, une plateforme numérique pour l'agriculture de
précision en Tunisie, combinant IoT (deux nœuds ESP32), vision par ordinateur (8 modèles
YOLO v11), un assistant conversationnel souverain en Darija (Labess-7B + RAG ChromaDB)
et un pipeline MLOps complet (DVC + MLflow). L'architecture repose sur un backend
FastAPI (50+ endpoints REST), un frontend React/Vite PWA offline, une base PostgreSQL
de 30+ tables et une détection d'anomalies télémétriques par IsolationForest
(Silhouette Score : 0,73). Le projet a été conduit en Scrum sur 6 sprints / 14 semaines
avec la méthodologie CRISP-DM pour la partie Data Science.

\medskip\noindent\textbf{Mots-clés :} Agriculture de précision, YOLO, IoT, ESP32, RAG, Darija, MLOps, FastAPI, Apiculture, Tunisie.

\clearpage

%% Abstract EN
\chapter*{Abstract}\thispagestyle{empty}
\addcontentsline{toc}{chapter}{Abstract}

\noindent\textbf{Title:} Smart Farm AI v3.0 — Multi-Tenant Agro-Technological Platform

\medskip
This report presents \farmAI{} v3.0, an innovative digital platform for precision
agriculture in Tunisia. The platform integrates IoT (two ESP32 nodes), computer vision
(8 YOLO v11 models covering 40 classes), a sovereign conversational assistant in Tunisian
Darija dialect (Labess-7B + RAG with ChromaDB), and a complete MLOps pipeline (DVC +
MLflow). The system features a FastAPI backend (50+ REST endpoints), a React/Vite
offline-capable PWA frontend, and IsolationForest-based telemetry anomaly detection
(silhouette score: 0.73). The project was conducted using Scrum (6 sprints, 14 weeks)
combined with CRISP-DM methodology.

\medskip\noindent\textbf{Keywords:} Precision agriculture, YOLO, IoT, ESP32, RAG, Darija, MLOps, FastAPI, Beekeeping, Tunisia.

\clearpage

%% Table des matières
\tableofcontents\clearpage
\listoffigures\addcontentsline{toc}{chapter}{Liste des Figures}\clearpage
\listoftables\addcontentsline{toc}{chapter}{Liste des Tableaux}\clearpage

%% Abréviations
\chapter*{Liste des Abréviations}\thispagestyle{empty}
\addcontentsline{toc}{chapter}{Liste des Abréviations}
\begin{longtable}{@{}lp{10cm}@{}}
\toprule\textbf{Abréviation}&\textbf{Signification}\\\midrule\endhead\bottomrule\endfoot
API     & Application Programming Interface \\
AVFA    & Agence de Vulgarisation et de Formation Agricoles \\
CRISP-DM& Cross-Industry Standard Process for Data Mining \\
DVC     & Data Version Control \\
ESP32   & Microcontrôleur Wi-Fi/Bluetooth (Espressif Systems) \\
FAO     & Food and Agriculture Organization (ONU) \\
IoT     & Internet of Things \\
JWT     & JSON Web Token \\
LLM     & Large Language Model \\
mAP     & Mean Average Precision \\
MLOps   & Machine Learning Operations \\
OBB     & Oriented Bounding Box \\
OTP     & One-Time Password \\
PWA     & Progressive Web App \\
RAG     & Retrieval-Augmented Generation \\
REST    & Representational State Transfer \\
SIG/GIS & Système d'Information Géographique \\
STT/TTS & Speech-to-Text / Text-to-Speech \\
UTAP    & Union Tunisienne de l'Agriculture et de la Pêche \\
YOLO    & You Only Look Once \\
\bottomrule
\end{longtable}
\clearpage

%% ──────────────────────────────────────────────────────────────
%%  CORPS DU RAPPORT (numérotation arabe)
%% ──────────────────────────────────────────────────────────────
\pagenumbering{arabic}\setcounter{page}{1}


% ============================================================
\chapter{Contexte g\'en\'eral et \'etude pr\'eliminaire}
% ============================================================

\section{Introduction}

Dans un contexte mondial marqu\'e par l\'acc\'el\'eration de la transformation num\'erique,
l\'intelligence artificielle s\'impose comme un levier strat\'egique pour moderniser des
secteurs traditionnels tels que l\'agriculture. La Tunisie, dont le secteur agricole
repr\'esente \textbf{10,3\% du PIB} et \textbf{14,1\% de l\'emploi national} (INS, 2023),
fait face \ des d\'efis structurels : stress hydrique, ravageurs responsables de
20--40\% des pertes, d\'eclin des pollinisateurs et vieillissement des exploitants.

Ce chapitre pose les fondations du projet \farmAI{} v3.0 en pr\'esentant l\'organisme
d\'accueil, le contexte du projet, la probl\'ematique identifi\'ee et les m\'ethodologies
de gestion de projet \'evalu\'ees.

\section{Pr\'esentation de l\'organisme d\'accueil}

\subsection{Mission et Activit\'es}

\begin{tcolorbox}[infobox,title={intech solution --- La Marsa}]
\textbf{intech solution} est une soci\'et\'e tunisienne de conseil et d\'ing\'enierie
informatique sp\'ecialis\'ee dans le d\'eveloppement de solutions num\'eriques sur mesure.
Implant\'ee \ La Marsa (Tunis), elle accompagne des entreprises et institutions dans
leur transformation digitale en proposant des services couvrant le d\'eveloppement
logiciel, l\'intelligence artificielle, l\'IoT et l\'int\'egration de syst\emes.
\end{tcolorbox}

Ses missions principales incluent :
\begin{itemize}
  \item Conception et d\'eveloppement d\'applications web et mobiles sur mesure.
  \item Int\'egration de solutions bas\'ees sur l\'intelligence artificielle et le Machine Learning.
  \item D\'eploiement de syst\emes IoT pour l\'industrie, l\'agriculture et les services.
  \item Conseil en architecture logicielle et transformation digitale.
  \item Formation et accompagnement des \'equipes clients.
\end{itemize}

\subsection{Domaines d\'Expertise et Points Forts}

\begin{table}[H]\centering
\caption{Domaines d\'expertise --- intech solution}
\label{tab:intech_expertise}
\begin{tabular}{ll}
\toprule
\textbf{Domaine} & \textbf{Technologies / Solutions} \\
\midrule
D\'eveloppement Web      & React, Vue.js, FastAPI, Django, Node.js \\
Intelligence Artificielle & YOLO, scikit-learn, TensorFlow, LLM, RAG \\
Internet des Objets      & ESP32, Arduino, MQTT, PlatformIO \\
Bases de donn\'ees        & PostgreSQL, MongoDB, ChromaDB, Redis \\
DevOps \& MLOps           & Docker, DVC, MLflow, GitHub Actions \\
Cloud                    & Oracle Cloud, AWS, Azure \\
\bottomrule
\end{tabular}
\end{table}

Les points forts d\'intech solution r\'esident dans sa capacit\'e \ allier expertise
technique de pointe et connaissance approfondie du tissu \'economique local tunisien,
lui permettant de proposer des solutions adapt\'ees aux contraintes sp\'ecifiques du march\'e
national (connectivit\'e, co\^uts, langues locales).

\subsection{Vision et Objectifs \ Long Terme}

La vision strat\'egique d\'intech solution s\'articule autour de trois axes :

\begin{enumerate}
  \item \textbf{Souverainet\'e technologique :} D\'evelopper des solutions num\'eriques
        enti\erement ma\^itris\'ees, r\'eduisant la d\'ependance aux plateformes cloud
        \'etrang\eres et aux mod\eles propri\'etaires.
  \item \textbf{Impact local :} Contribuer \ la transformation num\'erique des secteurs
        cl\'es de l\'\'economie tunisienne --- agriculture, sant\'e, \'education --- avec des
        outils adapt\'es aux r\'ealit\'es locales.
  \item \textbf{Innovation ouverte :} Encourager l\'utilisation de technologies open-source
        et contribuer \ l\'\'ecosyst\eme technologique tunisien par la formation et le
        partage de connaissances.
\end{enumerate}

\subsection{Informations cl\'es sur intech solution La Marsa}

\begin{table}[H]\centering
\caption{Fiche signal\'etique --- intech solution}
\label{tab:intech_fiche}
\begin{tabular}{ll}
\toprule
\textbf{Champ}            & \textbf{Information} \\
\midrule
Raison sociale            & intech solution \\
Si\ege social             & La Marsa, Tunis, Tunisie \\
Secteur d\'activit\'e       & Informatique \& Conseil en Technologies \\
Sp\'ecialit\'es             & IA, IoT, D\'eveloppement logiciel, MLOps \\
Cadre du stage            & Projet de fin de cycle d\'Ing\'enierie \\
Dur\'ee du stage            & 14 semaines (f\'evrier -- mai 2025) \\
Encadrant entreprise      & [Encadrant intech solution] \\
\bottomrule
\end{tabular}
\end{table}

\subsection{Structure organisationnelle}

\begin{figure}[H]\centering
\begin{tikzpicture}[
  node distance=1.4cm and 2.2cm,
  box/.style={rectangle, rounded corners=4pt, draw=#1!70!black, fill=#1!12,
    text width=3cm, minimum height=1cm, align=center, font=\small\bfseries},
  arr/.style={->, >=Stealth, thick, farmgray!70}
]
\node[box=farmblue, text width=4cm] (dg) {Direction G\'en\'erale};
\node[box=farmgreen, below left=of dg]  (dev) {D\'epartement\\ D\'eveloppement};
\node[box=farmgreen, below=of dg]       (ia)  {D\'epartement\\ IA \& Data};
\node[box=farmgreen, below right=of dg] (iot) {D\'epartement\\ IoT \& Syst\emes};
\node[box=farmorange, below=of dev]     (web) {Web \&\\ Mobile};
\node[box=farmorange, below=of ia]      (ml)  {Machine\\ Learning};
\node[box=farmorange, below=of iot]     (emb) {Syst\emes\\ Embarqu\'es};
\draw[arr] (dg) -- (dev);
\draw[arr] (dg) -- (ia);
\draw[arr] (dg) -- (iot);
\draw[arr] (dev) -- (web);
\draw[arr] (ia)  -- (ml);
\draw[arr] (iot) -- (emb);
\end{tikzpicture}
\caption{Structure organisationnelle --- intech solution}
\label{fig:org_chart}
\end{figure}

\section{Contexte du projet}

\subsection{\'Etude de l\'existant}

L\'\'etude de l\'existant a permis d\'\'evaluer les solutions agritech disponibles sur le
march\'e tunisien et international :

\begin{table}[H]\centering
\caption{Benchmark agritech : \farmAI{} vs solutions existantes}
\label{tab:benchmark}
\small
\begin{tabular}{lcccccc}
\toprule
\textbf{Crit\ere}
  &\textbf{Agrivi}&\textbf{FarmLogs}&\textbf{Trimble}
  &\textbf{AgroMon.}&\textbf{Local}&\textbf{\farmAI}\\
\midrule
IoT ESP32 int\'egr\'e        &\no &\no &\yes&\no &\no &\yes\\
YOLO vision pathologies  &\no &\no &\no &\no &\no &\yes\\
LLM Darija tunisien      &\no &\no &\no &\no &\no &\yes\\
RAG documents locaux     &\no &\no &\no &\no &\no &\yes\\
Module apiculture complet&\no &\no &\no &\no &\no &\yes\\
MLOps DVC + MLflow       &\no &\no &\yes&\no &\no &\yes\\
PWA offline Worker       &\no &\yes&\no &\no &\no &\yes\\
Open-source / gratuit    &\no &\no &\no &\no &\yes&\yes\\
\midrule
\textbf{Score total}&0/8&1/8&2/8&0/8&2/8&\textbf{8/8}\\
\bottomrule
\end{tabular}
\end{table}

\noindent Ce benchmark r\'ev\ele qu\'aucune solution existante ne combine IoT, vision IA,
assistant en dialecte local et MLOps en un syst\eme int\'egr\'e accessible aux agriculteurs
tunisiens.

\subsection{Probl\'ematique}

\begin{tcolorbox}[warningbox,title={Probl\'ematique}]
\itshape Comment concevoir une plateforme IA souveraine et accessible pour les
agriculteurs tunisiens, combinant surveillance IoT temps r\'eel, d\'etection de pathologies
par vision par ordinateur et assistance conversationnelle en dialecte local (Darija),
tout en garantissant la reproductibilit\'e des mod\eles via des pratiques MLOps
industrielles et un d\'eploiement \'economiquement viable ?
\end{tcolorbox}

\subsection{Objectifs}

Le projet \farmAI{} v3.0 vise les objectifs suivants :

\begin{enumerate}
  \item \textbf{Surveillance IoT temps r\'eel :} Collecter et analyser les donn\'ees
        t\'el\'em\'etriques des deux n\oe{}uds ESP32 (irrigation et rucher) avec d\'etection
        automatique d\'anomalies.
  \item \textbf{Vision par ordinateur :} D\'eployer 8 mod\eles YOLO v11 couvrant 40 classes
        de pathologies agricoles et apicoles avec une latence $< ms sur CPU.
  \item \textbf{Assistant souverain :} Fournir un assistant conversationnel en Darija
        tunisien bas\'e sur Labess-7B et RAG ChromaDB, op\'erationnel sans connexion cloud.
  \item \textbf{MLOps industriel :} Impl\'ementer un pipeline DVC + MLflow garantissant
        la reproductibilit\'e et le versionnement des mod\eles.
  \item \textbf{Accessibilit\'e :} D\'eployer sur Oracle Cloud Free Tier avec un mode
        d\'egrad\'e (LITE\_MODE) pour les contraintes mat\'erielles.
\end{enumerate}

\subsection{Solution propos\'ee}

\begin{tcolorbox}[infobox,title={Les 4 piliers de Smart Farm AI v3.0}]
\begin{enumerate}
  \item \textbf{IoT Temps R\'eel :} NODE\_A (irrigation) + NODE\_B (rucher) ---
        HTTP POST toutes les 10~s, d\'etection anomalies IsolationForest.
  \item \textbf{Vision IA :} 8 mod\eles YOLO v11, 40 classes, inférence locale $< ms,
        scanner PWA pour les ouvriers.
  \item \textbf{Assistant Darija :} Labess-7B (Ollama) + RAG ChromaDB (corpus UTAP/AVFA),
        STT/TTS, fallback Groq.
  \item \textbf{MLOps :} DVC + MLflow --- reproductibilit\'e, registre de mod\eles,
        alertes anomalies WebSocket.
\end{enumerate}
\end{tcolorbox}

L\'architecture retenue est une plateforme web multi-tenant reposant sur un backend
\textbf{FastAPI} (50+ endpoints REST), un frontend \textbf{React/Vite PWA} capable
de fonctionner hors ligne, une base \textbf{PostgreSQL} de 30+ tables, et une
infrastructure \textbf{Docker Compose} d\'eployable sur Oracle Cloud ARM64.

\subsection{Motivations}

Les motivations ayant guid\'e ce projet sont d\'ordre technique, \'economique et social :

\begin{description}
  \item[\textbf{Motivation technique :}] Appliquer des technologies de pointe
        (YOLO v11, LLM, RAG, IoT, MLOps) dans un contexte r\'eel et contraignant,
        sans GPU et avec des ressources cloud gratuites.
  \item[\textbf{Motivation \'economique :}] L\'agriculture tunisienne perd 20--40\%
        de sa production \ cause des ravageurs et maladies. Une d\'etection pr\'ecoce
        par IA peut r\'eduire ces pertes de mani\ere significative.
  \item[\textbf{Motivation sociale :}] Proposer un outil accessible en Darija tunisien,
        adapt\'e aux r\'ealit\'es linguistiques et culturelles des agriculteurs locaux,
        souvent peu familiers des interfaces en fran\c{c}ais ou en anglais.
  \item[\textbf{Motivation acad\'emique :}] Valider les comp\'etences acquises en
        Data Science et Intelligence Artificielle \ travers un projet int\'egrateur
        couvrant l\'ensemble du cycle de vie d\'un syst\eme IA (donn\'ees, mod\'elisation,
        d\'eploiement, monitoring).
\end{description}

\subsection{Contraintes}

\begin{table}[H]\centering
\caption{Contraintes du projet \farmAI{} v3.0}
\label{tab:contraintes}
\begin{tabular}{lp{4.5cm}p{5cm}}
\toprule
\textbf{Type} & \textbf{Contrainte} & \textbf{Solution retenue} \\
\midrule
Mat\'erielle  & Pas de GPU disponible & Inférence CPU YOLOv11 ($< ms), LITE\_MODE \\
Cloud        & Oracle Free Tier ARM64 & LITE\_MODE d\'esactive YOLO et LLM si n\'ecessaire \\
Linguistique & Dialecte Darija non standardis\'e & Labess-7B natif tunisien + corpus UTAP/AVFA \\
Temporelle   & 14 semaines (6 sprints) & Planification Scrum stricte, jalons hebdomadaires \\
Budget       & Z\'ero co\^ut infrastructure & Technologies 100\% open-source et cloud gratuit \\
S\'ecurit\'e   & Donn\'ees agricoles sensibles & JWT + OTP, isolation multi-tenant, bcrypt \\
\bottomrule
\end{tabular}
\end{table}

\section{Gestion de projet}

\subsection{M\'ethodologie}

La gestion d\'un projet m\^elant d\'eveloppement logiciel agile et pipeline Data Science
n\'ecessite l\'articulation de deux types de m\'ethodologies compl\'ementaires :
\begin{itemize}
  \item \textbf{M\'ethodologies agiles} pour le d\'eveloppement logiciel it\'eratif (Scrum).
  \item \textbf{M\'ethodologies Data Science} pour le cycle de vie des mod\eles
        (KDD, SEMMA, CRISP-DM).
\end{itemize}

\subsection{Knowledge Discovery in Databases (KDD)}

Le processus KDD (Fayyad et al., 1996) est le cadre fondateur de la d\'ecouverte de
connaissances dans les bases de donn\'ees. Il d\'efinit un processus it\'eratif en
\textbf{9 \'etapes} :

\begin{table}[H]\centering
\caption{Processus KDD --- 9 \'etapes}
\label{tab:kdd}
\small
\begin{tabular}{clp{7cm}}
\toprule
\textbf{\'Etape} & \textbf{Phase} & \textbf{Description} \\
\midrule
1 & S\'election         & Identification des donn\'ees pertinentes (t\'el\'em\'etrie IoT, images) \\
2 & Pr\'e-traitement   & Nettoyage, gestion des valeurs manquantes \\
3 & Transformation     & Normalisation, augmentation, extraction de features \\
4 & Fouille            & Application de YOLO, IsolationForest, ChromaDB \\
5 & Interpr\'etation   & mAP50, Silhouette Score, alertes \\
6--9 & It\'eration      & R\'eentra\^inement, ajustement, d\'eploiement \\
\bottomrule
\end{tabular}
\end{table}

\noindent\textbf{Limite :} KDD est une m\'ethodologie g\'en\'erique peu prescriptive sur
les aspects de d\'eploiement et de monitoring continu, ce qui le rend insuffisant pour
un projet MLOps industriel.

\subsection{SEMMA}

SEMMA (SAS Institute, 1990s) est une m\'ethodologie orient\'ee donn\'ees structur\'ee en
5 phases : \textbf{S}ample, \textbf{E}xplore, \textbf{M}odify, \textbf{M}odel,
\textbf{A}ssess.

\begin{table}[H]\centering
\caption{M\'ethodologie SEMMA appliqu\'ee \ \farmAI{}}
\label{tab:semma}
\begin{tabular}{clp{7.5cm}}
\toprule
\textbf{Phase} & \textbf{Nom} & \textbf{Application \farmAI{}} \\
\midrule
S & Sample  & \'Echantillonnage t\'el\'em\'etrie ESP32 (CSV) et images Roboflow \\
E & Explore & EDA : distributions IoT, d\'es\'equilibre classes YOLO \\
M & Modify  & Normalisation, augmentation Mosaic, split 80/10/10 \\
M & Model   & YOLO v11, IsolationForest, Labess-7B \\
A & Assess  & mAP50 $\geq$ 0.85, Silhouette $\geq$ 0.70 \\
\bottomrule
\end{tabular}
\end{table}

\noindent\textbf{Limite :} SEMMA est centr\'ee sur SAS et ne couvre pas le versionnement
des mod\eles, le monitoring en production ni les pratiques MLOps modernes.

\subsection{CRISP-DM}

CRISP-DM (Cross-Industry Standard Process for Data Mining) est la m\'ethodologie
la plus adopt\'ee en industrie (60\% des projets Data Science selon KDnuggets, 2021).
Elle d\'efinit \textbf{6 phases cycliques et it\'eratives} :

\begin{figure}[H]\centering
\begin{tikzpicture}[x=1cm,y=1cm]
\tikzset{
  ph/.style={circle,draw=#1!80!black,fill=#1!20,text width=2cm,align=center,
    font=\tiny\bfseries,inner sep=2pt,minimum size=2.2cm},
}
\def\R{3.8}
\node[ph=farmblue]  (biz)   at (0*60:\R)   {Compreh.\\ M\'etier};
\node[ph=farmgreen] (data)  at (-1*60:\R)  {Compreh.\\ Donn\'ees};
\node[ph=farmorange](prep)  at (-2*60:\R)  {Pr\'ep.\\ Donn\'ees};
\node[ph=farmred]   (model) at (180:\R)    {Mod\'eli-\\ sation};
\node[ph=farmblue!70](eval) at (3*60:\R)   {\'Evalua-\\ tion};
\node[ph=farmgreen!70](dep) at (2*60:\R)   {D\'eploie-\\ ment};
\foreach \a/\b in {biz/data,data/prep,prep/model,model/eval,eval/dep,dep/biz}
  \draw[->,>=Stealth,thick,farmgray!60,bend left=18](\a)to(\b);
\draw[->,>=Stealth,thick,farmred!70,dashed,bend left=50]
  (eval)to node[right,font=\tiny\itshape,pos=0.5]{mAP$<.70}(model);
\node[ellipse,draw=farmgray!30,fill=white,text width=2.2cm,align=center,
  font=\small\bfseries\color{farmblue},inner sep=4pt]at(0,0){\farmAI\\Data\\ Pipeline};
\end{tikzpicture}
\caption{Cycle CRISP-DM --- \farmAI{} v3.0}
\label{fig:crisp_dm_ch1}
\end{figure}

\begin{table}[H]\centering
\caption{CRISP-DM appliqu\'e \ \farmAI{} v3.0}
\label{tab:crisp_ch1}
\begin{tabular}{clp{8cm}}
\toprule
\textbf{Ph.} & \textbf{Phase} & \textbf{Actions dans \farmAI{}} \\
\midrule
1 & Compr\'ehension m\'etier  & Pathologies cibles, KPIs (mAP50$\geq.85), contraintes ($< ms) \\
2 & Compr\'ehension donn\'ees & Datasets Roboflow YOLO, CSV IoT ESP32, corpus UTAP/AVFA \\
3 & Pr\'eparation donn\'ees   & DVC \texttt{build\_dataset} : normalisation, augmentation, split 80/10/10 \\
4 & Mod\'elisation            & 8 YOLO v11, IsolationForest (n=150), ChromaDB HNSW \\
5 & \'Evaluation              & mAP50=91.4\%, Silhouette=0.73 --- it\'eration si mAP50$<.70 \\
6 & D\'eploiement             & MLflow Registry, Docker Compose, Oracle Cloud LITE\_MODE \\
\bottomrule
\end{tabular}
\end{table}

\section{Choix de la m\'ethodologie : CRISP-DM}

Apr\es comparaison des trois m\'ethodologies, CRISP-DM a \'et\'e retenu comme cadre
m\'ethodologique principal du projet \farmAI{} v3.0 pour les raisons suivantes :

\begin{table}[H]\centering
\caption{Comparaison KDD / SEMMA / CRISP-DM}
\label{tab:methodo_compare}
\small
\begin{tabular}{lccc}
\toprule
\textbf{Crit\ere}              & \textbf{KDD} & \textbf{SEMMA} & \textbf{CRISP-DM} \\
\midrule
Ind\'ependance outil             & \yes         & \no (SAS)      & \yes \\
Couverture d\'eploiement         & Partielle    & Non            & \yes \\
Crit\ere d\'arr\^et explicite   & \no          & \no            & \yes \\
Adopt\'e en industrie            & Moyen        & Faible         & \yes (60\%) \\
Compatible MLOps                 & Partiel      & \no            & \yes \\
It\'erativit\'e sur les mod\eles & Oui          & Oui            & \yes \\
\midrule
\textbf{D\'ecision}             & --           & --             & \textbf{\yes~Choisi} \\
\bottomrule
\end{tabular}
\end{table}

\begin{tcolorbox}[infobox,title={Pourquoi CRISP-DM pour \farmAI{} ?}]
\begin{itemize}
  \item \textbf{Crit\ere d\'arr\^et quantitatif :} La boucle d\'it\'eration
        (mAP50 $\geq$ 0.70) \'evite la sur-ing\'enierie.
  \item \textbf{Phase D\'eploiement explicite :} Aligne avec le pipeline
        MLflow Registry + Docker Compose.
  \item \textbf{Ind\'ependance technologique :} Compatible avec
        DVC, MLflow, scikit-learn --- aucune d\'ependance \ un outil propri\'etaire.
  \item \textbf{Compl\'ementarit\'e Scrum :} Les phases CRISP-DM s\'int\egrent
        naturellement dans les sprints S4 (Vision IA) et S6 (MLOps).
\end{itemize}
\end{tcolorbox}

\section{Conclusion}

Ce chapitre a pos\'e les bases du projet \farmAI{} v3.0 en pr\'esentant l\'organisme
d\'accueil intech solution (La Marsa), le contexte agricole tunisien, et la
probl\'ematique identifi\'ee. L\'\'etude de l\'existant a confirm\'e l\'absence de solution
int\'egrant IoT, vision IA et assistant en Darija adapt\'e au march\'e local.

La m\'ethodologie \textbf{CRISP-DM} a \'et\'e choisie pour sa rigueur, sa couverture du
cycle complet (de la compr\'ehension m\'etier au d\'eploiement) et sa compatibilit\'e avec
les pratiques MLOps modernes. Le chapitre suivant pr\'esentera l\'architecture technique
d\'etaill\'ee et les choix de conception du syst\eme.


% ============================================================
\chapter{Méthodologies de Conduite de Projet}
% ============================================================

\noindent Ce chapitre présente les deux méthodologies structurant la conduite du projet :
\textbf{Scrum} pour la gestion agile du développement logiciel en 6 sprints, et
\textbf{CRISP-DM} pour la rigueur scientifique du cycle Data Science. Ces deux
approches sont complémentaires : Scrum cadence les livraisons, CRISP-DM guide
les itérations de modélisation.

\section{Scrum — 6 Sprints}

\begin{table}[H]\centering
\caption{Planification des 6 Sprints — Smart Farm AI}
\label{tab:sprints}
\small
\begin{tabular}{cccp{5.5cm}l}
\toprule
\textbf{Spr.}&\textbf{Dur.}&\textbf{Titre}&\textbf{Livrables}&\textbf{Jalon}\\
\midrule
S1&S0--S2&Infra \& Auth&FastAPI, JWT+OTP, PostgreSQL 30+ tables, React scaffold&Auth OK\\
S2&S2--S4&IoT \& Télémétrie&NODE\_A+B firmware HTTP POST, Wokwi sim., dashboard Recharts&IoT Live\\
S3&S4--S7&Apiculture&HiveIntel 7 sous-modules, CRUD ruchers/ruches/visites&HiveIntel v1.0\\
S4&S7--S10&Vision IA&8 YOLO v11, Scanner PWA 13 modèles, Map SIG Leaflet&YOLO Prod\\
S5&S10--S12&Assistant&Labess-7B RAG Darija, STT/TTS, LLaVA, Lite Mode Oracle&RAG Darija\\
S6&S12--S14&MLOps&DVC+MLflow, IsolationForest, Docker Compose, Oracle Cloud&Soutenance\\
\bottomrule
\end{tabular}
\end{table}

\subsection{Diagramme Gantt — 14 semaines}

\begin{figure}[H]\centering
\begin{tikzpicture}[x=0.85cm,y=1cm]
\tikzset{
  spr/.style={fill=#1,draw=#1!60!black,rounded corners=2pt,opacity=0.88},
  ms/.style={diamond,draw=farmgreen!80!black,fill=farmgreen!30,minimum size=0.5cm,inner sep=0pt},
  sout/.style={diamond,draw=red!70!black,fill=red!25,minimum size=0.6cm,inner sep=0pt}
}
\draw[->,thick,gray](-0.3,0)--(15.2,0)node[right,font=\small]{Semaines};
\foreach \w in {0,...,14}{
  \draw[gray!40,thin](\w,0.05)--(\w,-0.15);
  \node[below,font=\tiny\color{gray}]at(\w,-0.18){S\w};
}
\filldraw[spr=farmblue!50] (0,0.15)rectangle(2,0.75);
\node[font=\tiny\bfseries\color{white}]at(1,0.45){S1 Auth};
\filldraw[spr=farmblue!65] (2,0.15)rectangle(4,0.75);
\node[font=\tiny\bfseries\color{white}]at(3,0.45){S2 IoT};
\filldraw[spr=farmblue!75] (4,0.15)rectangle(7,0.75);
\node[font=\tiny\bfseries\color{white}]at(5.5,0.45){S3 Apiculture};
\filldraw[spr=farmblue!85] (7,0.15)rectangle(10,0.75);
\node[font=\tiny\bfseries\color{white}]at(8.5,0.45){S4 Vision IA};
\filldraw[spr=farmblue!90] (10,0.15)rectangle(12,0.75);
\node[font=\tiny\bfseries\color{white}]at(11,0.45){S5 Assistant};
\filldraw[spr=farmblue!95] (12,0.15)rectangle(14,0.75);
\node[font=\tiny\bfseries\color{white}]at(13,0.45){S6 MLOps};
\foreach \x/\lbl in {2/{Auth\\JWT},4/{IoT Live\\Dash.},7/{HiveIntel\\v1.0},10/{YOLO\\Prod},12/{RAG\\Darija}}{
  \node[ms](m)at(\x,-0.7){};
  \node[below=0.04cm of m,font=\tiny,align=center]{\lbl};
  \draw[dashed,gray!50,thin](\x,0)--(\x,-0.6);
}
\node[sout](ms)at(14,-0.7){};
\node[below=0.04cm of ms,font=\tiny\bfseries,text=red!70!black,align=center]{Souten.\\v3.0};
\draw[dashed,red!50,thin](14,0)--(14,-0.6);
\node[ms,scale=0.75]at(0,-1.6){};
\node[right,font=\tiny]at(0.3,-1.6){Jalon sprint};
\node[sout,scale=0.75]at(3,-1.6){};
\node[right,font=\tiny]at(3.3,-1.6){Soutenance finale};
\end{tikzpicture}
\caption{Diagramme Gantt — Planification Scrum (6 Sprints, 14 semaines)}
\label{fig:gantt}
\end{figure}

\section{CRISP-DM}

\begin{table}[H]\centering
\caption{CRISP-DM appliqué à Smart Farm AI v3.0}
\label{tab:crisp}
\begin{tabular}{clp{9cm}}
\toprule
\textbf{Ph.}&\textbf{Phase}&\textbf{Actions dans FARM AI}\\
\midrule
1&Compréhension métier&Pathologies cibles, KPIs (mAP50$\geq$0.85), contraintes ($<$200 ms)\\
2&Compréhension données&Datasets Roboflow YOLO, CSV IoT ESP32, corpus UTAP/AVFA (RAG)\\
3&Préparation données&DVC \texttt{build\_dataset}: normalisation CSV, augmentation, split 80/10/10\\
4&Modélisation&8 YOLO v11, IsolationForest (n=150), ChromaDB HNSW, prompts Labess-7B\\
5&Évaluation&mAP50=91.4\%, Silhouette=0.73 — itération si mAP50$<$0.70\\
6&Déploiement&MLflow Registry, Docker Compose, Oracle Cloud LITE\_MODE\\
\bottomrule
\end{tabular}
\end{table}

\begin{figure}[H]\centering
\begin{tikzpicture}[x=1cm,y=1cm]
\tikzset{
  ph/.style={circle,draw=#1!80!black,fill=#1!20,text width=2cm,align=center,
    font=\tiny\bfseries,inner sep=2pt,minimum size=2.2cm},
  ann/.style={rectangle,rounded corners=2pt,draw=farmgray!50,fill=yellow!12,
    text width=2.8cm,align=left,font=\tiny,inner sep=3pt}
}
\def\R{3.8}
\node[ph=farmblue]  (biz)   at (0*60:\R)   {Compreh.\\ Métier};
\node[ph=farmgreen] (data)  at (-1*60:\R)  {Compreh.\\ Données};
\node[ph=farmorange](prep)  at (-2*60:\R)  {Prép.\\ Données};
\node[ph=farmred]   (model) at (180:\R)    {Modéli-\\ sation};
\node[ph=farmblue!70](eval) at (3*60:\R)   {Évalua-\\ tion};
\node[ph=farmgreen!70](dep) at (2*60:\R)   {Déploie-\\ ment};
\foreach \a/\b in {biz/data,data/prep,prep/model,model/eval,eval/dep,dep/biz}
  \draw[->,>=Stealth,thick,farmgray!60,bend left=18](\a)to(\b);
\draw[->,>=Stealth,thick,farmred!70,dashed,bend left=50]
  (eval)to node[right,font=\tiny\itshape,pos=0.5]{mAP$<$0.70}(model);
\node[ellipse,draw=farmgray!30,fill=white,text width=2.2cm,align=center,
  font=\small\bfseries\color{farmblue},inner sep=4pt]at(0,0){\farmAI\\Data\\ Pipeline};
\end{tikzpicture}
\caption{Cycle CRISP-DM — Smart Farm AI v3.0}
\label{fig:crisp_dm}
\end{figure}

\medskip\noindent\textbf{Conclusion du chapitre :} Les méthodologies Scrum et CRISP-DM
ont été choisies pour leur complémentarité : Scrum permet des livraisons fonctionnelles
toutes les 2 à 3 semaines avec adaptation aux imprévus, tandis que CRISP-DM garantit
la rigueur du processus Data Science avec un critère d'arrêt explicite (mAP50 $\geq$ 0.70).
Le chapitre suivant traduit ces choix méthodologiques en décisions d'architecture.

% ============================================================
\chapter{Architecture et Conception}
% ============================================================

\noindent Ce chapitre détaille l'architecture logicielle et matérielle de \farmAI{} v3.0.
Après les diagrammes de cas d'utilisation UML, les diagrammes de séquence illustrent
les flux d'interaction clés. L'architecture en 5 couches, la Clean Architecture
backend et la sécurité du système complètent la vue de conception.

\section{Diagramme de cas d'utilisation global}

\begin{figure}[H]\centering
\begin{tikzpicture}[x=1cm,y=1cm,
  uc/.style={ellipse,draw=farmblue!70,fill=farmblue!8,text width=2.6cm,align=center,
    font=\tiny,inner sep=2pt,minimum height=0.7cm},
  incl/.style={dashed,->,>=Stealth,farmblue!60,font=\tiny},
  ext/.style={dashed,->,>=Stealth,farmorange!70,font=\tiny},
  assoc/.style={-,thick,farmgray!60}]

\draw[rounded corners=10pt,draw=farmblue,fill=farmblue!3,line width=1.2pt]
  (1.5,-9.5)rectangle(12.5,1.0);
\node[font=\normalsize\bfseries\color{farmblue},above]at(7,1.0){\farmAI{} v3.0};

\stickactor{-0.5}{-0.5}{Owner\\(Propriétaire)}
\coordinate(owner)at(-0.5,-1.5);
\stickactor{-0.5}{-7.0}{Worker\\(Ouvrier)}
\coordinate(worker)at(-0.5,-8.0);
\stickactor{14.5}{-0.5}{ESP32\\NODE\_A/B}
\coordinate(esp32)at(14.5,-1.5);
\stickactor{14.5}{-7.0}{Ollama\\Labess-7B}
\coordinate(ollama)at(14.5,-8.0);

\node[uc](uc_auth)  at(3.5, 0)    {S'authentifier\\JWT + OTP};
\node[uc](uc_farm)  at(3.5,-1.8)  {Gérer\\Fermes};
\node[uc](uc_dash)  at(3.5,-3.5)  {Tableau\\de Bord};
\node[uc](uc_tele)  at(3.5,-5.2)  {Télémétrie\\IoT};
\node[uc](uc_alert) at(3.5,-6.9)  {Alertes\\IoT};
\node[uc](uc_map)   at(3.5,-8.5)  {Map SIG\\Leaflet};

\node[uc](uc_scan)  at(7.5,-1.5)  {Scanner\\YOLO (13)};
\node[uc](uc_hive)  at(7.5,-4.5)  {Gestion\\Apiculture};
\node[uc](uc_work)  at(7.5,-7.5)  {Gestion\\Ouvriers};

\node[uc](uc_chat)  at(11,-1.5)   {Assistant\\Darija};
\node[uc](uc_stt)   at(11,-3.2)   {STT / TTS};
\node[uc](uc_iot)   at(11,-5.5)   {Envoi\\Télémétrie};
\node[uc](uc_anom)  at(11,-7.5)   {Détection\\Anomalies};

\foreach \u in {uc_auth,uc_farm,uc_dash,uc_tele,uc_alert,uc_map,uc_scan,uc_hive,uc_work,uc_chat}
  \draw[assoc](owner)--(\u);
\foreach \u in {uc_auth,uc_scan,uc_work}
  \draw[assoc](worker)--(\u);
\draw[assoc](esp32)--(uc_iot);
\draw[assoc](ollama)--(uc_chat);
\draw[incl](uc_chat)--node[above,sloped]{\texttt{<<include>>}}(uc_stt);
\draw[incl](uc_dash)--node[above,sloped,pos=0.4]{\texttt{<<include>>}}(uc_tele);
\draw[ext] (uc_iot)--node[right]{\texttt{<<extend>>}}(uc_anom);
\end{tikzpicture}
\caption{Diagramme de cas d'utilisation global}
\label{fig:uc_global}
\end{figure}

\section{Architecture système en 5 couches}

\begin{figure}[H]\centering
\begin{tikzpicture}[x=1cm,y=1cm,
  comp/.style={rectangle,rounded corners=3pt,draw=#1!60!black,fill=#1!18,
    text width=2.8cm,minimum height=1.1cm,align=center,font=\tiny\bfseries,inner sep=3pt}]

\foreach \y/\col/\lbl/\desc in {
  0/farmblue/{Couche Clients}/{Owner Dashboard | Worker PWA | Map Center | Assistant},
  -2.6/farmgreen/{Couche API}/{JWT Middleware | FastAPI 50+ Endpoints | WebSocket | CORS},
  -5.2/farmorange/{Couche Services}/{CV YOLO | RAG ChromaDB | LLM Agent | Anomaly | Weather},
  -7.8/farmred/{Couche Données}/{PostgreSQL | ChromaDB | MLflow | DVC | Ollama},
  -10.4/farmgray/{Couche IoT}/{NODE\_A Irrigation | Wokwi Sim. | NODE\_B Rucher | PlatformIO}}{
  \draw[rounded corners=5pt,draw=\col,fill=\col!8,line width=1pt]
    (-8,\y-0.85)rectangle(8,\y+0.85);
  \node[font=\tiny\bfseries,text=\col!80!black]at(-6.5,\y+0.6){\lbl};
  \node[font=\tiny,text=farmgray,align=center]at(0,\y){\desc};
}
\foreach \y in {-0.85,-3.45,-6.05,-8.65}
  \draw[->,thick,>=Stealth,farmgray!60](0,\y)--(0,\y-0.8);
\node[right,font=\tiny\itshape,text=farmgray]at(0.2,-1.3){HTTPS/REST};
\node[right,font=\tiny\itshape,text=farmgray]at(0.2,-3.9){Appels Python};
\node[right,font=\tiny\itshape,text=farmgray]at(0.2,-6.5){SQLAlchemy};
\node[right,font=\tiny\itshape,text=farmgray]at(0.2,-9.1){HTTP POST};
\end{tikzpicture}
\caption{Architecture système en 5 couches — Smart Farm AI v3.0}
\label{fig:archi}
\end{figure}

\section{Clean Architecture du backend FastAPI}

\begin{figure}[H]\centering
\begin{tikzpicture}[x=1cm,y=1cm,
  lay/.style={rectangle,rounded corners=4pt,draw=#1!70!black,fill=#1!10,
    minimum width=13cm,minimum height=1.4cm,align=center,font=\small\bfseries\color{#1!80!black}},
  comp/.style={rectangle,rounded corners=2pt,draw=#1!60!black,fill=#1!20,
    text width=2.2cm,minimum height=1.0cm,align=center,font=\tiny\bfseries,inner sep=3pt}]
\foreach \y/\col/\items in {
  0/farmblue/{Routes API -- FastAPI Endpoints},
  -2.4/farmgreen/{Services -- Logique Métier},
  -4.8/farmorange/{Repositories -- Accès Données},
  -7.2/farmred/{Modèles -- SQLAlchemy ORM}}{
  \node[lay=\col]at(0,\y){};
  \node[font=\tiny\bfseries,text=\col!60]at(-5.8,\y+0.4){\items};
}
\foreach \x/\lbl in {-4.5/{Auth},  -1.5/{CV/YOLO}, 1.5/{IoT}, 4.5/{Apiculture}}
  \foreach \y in {0,-2.4,-4.8,-7.2}
    \node[comp=farmblue,minimum height=0.9cm,text width=1.8cm]at(\x,\y){\lbl};
\foreach \y in {-0.75,-3.15,-5.55}
  \draw[->,thick,>=Stealth,farmgray!70](0,\y)--(0,\y-0.8);
\node[rectangle,rounded corners=4pt,draw=farmgray!60,fill=farmgray!10,
  minimum width=13cm,minimum height=1.4cm,align=center,font=\small\bfseries\color{farmgray}]
  at(0,-9.6){Infrastructure : PostgreSQL | SQLite | ChromaDB | Ollama | YOLO .pt};
\draw[->,thick,>=Stealth,farmgray!70](0,-7.95)--(0,-8.85);
\end{tikzpicture}
\caption{Clean Architecture backend FastAPI}
\label{fig:clean}
\end{figure}

\section{Architecture IoT}

\begin{table}[H]\centering
\caption{Configuration des deux nœuds ESP32}
\label{tab:iot}
\begin{tabular}{llll}
\toprule
\textbf{Nœud}&\textbf{Capteur / Pin}&\textbf{Métrique}&\textbf{Seuil d'alerte}\\
\midrule
NODE\_A & Sol résistif GPIO34  & \code{soil}      & Sec$<$35\% / Humide$>$65\%\\
NODE\_A & Pression ADC GPIO35  & \code{pressure}  & Max 10 bar\\
NODE\_A & Débit YF-S201 GPIO32 & \code{flow}      & Fuite$<$0.5 L/min (pompe ON)\\
NODE\_A & DHT22 GPIO14         & \code{temp}      & Référence extérieure\\
\midrule
NODE\_B & Jauge poids GPIO34   & \code{weight}    & SWARM si $\Delta>$2 kg/60 min\\
NODE\_B & DHT22 GPIO14         & \code{ext\_temp/hum}& Hum$>$85\% WARN\\
NODE\_B & DS18B20 GPIO15       & \code{hive\_temp}& $<$30°C ou $>$38°C ALERT\\
\bottomrule
\end{tabular}
\end{table}

\begin{lstlisting}[style=arducstyle,caption={Envoi télémétrie HTTP POST — ESP32 NODE\_B (extrait)},label={lst:iot}]
void postMetric(const char* metric, float value) {
    HTTPClient http;
    String url = String("http://") + API_HOST + ":"
               + API_PORT + "/api/v1/iot/telemetry";
    http.begin(url);
    http.addHeader("Content-Type", "application/json");
    http.setTimeout(3000);
    String body = "{\"node\":\"NODE_B\","
                  "\"metric\":\"" + String(metric) + "\","
                  "\"value\":" + String(value, 4) + "}";
    int code = http.POST(body);  // HTTP 200 = success
    http.end();
}
\end{lstlisting}

\section{Diagrammes de séquence}

\subsection{Flux télémétrie IoT — ESP32 vers Dashboard}

\begin{figure}[H]\centering
\begin{tikzpicture}[x=1cm,y=1cm,
  actor/.style={rectangle,draw=farmblue!70,fill=farmblue!10,minimum width=2cm,
    minimum height=0.6cm,align=center,font=\tiny\bfseries},
  msg/.style={->,>=Stealth,thick,farmgray!80},
  ret/.style={->,>=Stealth,dashed,farmgray!50},
  lifeline/.style={dashed,farmgray!40,thin}
]
% Acteurs (colonnes)
\foreach \x/\lbl in {0/{ESP32\\NODE\_B}, 3.5/{FastAPI\\Backend}, 7/{PostgreSQL}, 10.5/{IsolationForest}, 14/{Dashboard\\WebSocket}} {
  \node[actor] at (\x, 0) {\lbl};
  \draw[lifeline] (\x,-0.3)--(\x,-9.5);
}
% Messages numérotés
\draw[msg](0,-1.0)--(3.5,-1.0)
  node[midway,above,font=\tiny]{1: POST /iot/telemetry \{node,metric,value\}};
\draw[msg](3.5,-2.0)--(7,-2.0)
  node[midway,above,font=\tiny]{2: INSERT TelemetryRecord};
\draw[ret](7,-2.7)--(3.5,-2.7)
  node[midway,below,font=\tiny]{3: id=42, OK};
\draw[msg](3.5,-3.7)--(10.5,-3.7)
  node[midway,above,font=\tiny]{4: score\_samples(features)};
\draw[ret](10.5,-4.4)--(3.5,-4.4)
  node[midway,below,font=\tiny]{5: score=0.12 (anomalie!)};
\draw[msg](3.5,-5.4)--(7,-5.4)
  node[midway,above,font=\tiny]{6: UPDATE is\_anomaly=True};
\draw[msg](3.5,-6.4)--(14,-6.4)
  node[midway,above,font=\tiny]{7: WS broadcast \{alert:"THERMAL\_FAULT"\}};
\draw[ret](3.5,-7.4)--(0,-7.4)
  node[midway,below,font=\tiny]{8: HTTP 200 OK};
% Boîtes d'activation
\foreach \x in {0,3.5,7,10.5,14}
  \filldraw[fill=farmblue!20,draw=farmblue!60,rounded corners=1pt]
    (\x-0.15,-0.3) rectangle (\x+0.15,-8.2);
\end{tikzpicture}
\caption{Diagramme de séquence — Flux télémétrie IoT (NODE\_B → Dashboard)}
\label{fig:seq_iot}
\end{figure}

\subsection{Flux authentification JWT + OTP}

\begin{figure}[H]\centering
\begin{tikzpicture}[x=1cm,y=1cm,
  actor/.style={rectangle,draw=farmblue!70,fill=farmblue!10,minimum width=2.2cm,
    minimum height=0.6cm,align=center,font=\tiny\bfseries},
  msg/.style={->,>=Stealth,thick,farmgray!80},
  ret/.style={->,>=Stealth,dashed,farmgray!50},
  lifeline/.style={dashed,farmgray!40,thin},
  note/.style={rectangle,rounded corners=2pt,draw=farmorange!50,fill=farmorange!10,
    font=\tiny,inner sep=3pt}
]
\foreach \x/\lbl in {0/{Client\\Navigateur}, 4/{FastAPI\\Backend}, 8/{PostgreSQL}, 12/{Service\\Email}} {
  \node[actor] at (\x, 0) {\lbl};
  \draw[lifeline] (\x,-0.3)--(\x,-11.5);
}
% Étape 1 : Inscription
\node[note]at(6,0.8){Phase 1 : Inscription};
\draw[msg](0,-1.0)--(4,-1.0)  node[midway,above,font=\tiny]{1: POST /auth/register \{email, password\}};
\draw[msg](4,-1.8)--(8,-1.8)  node[midway,above,font=\tiny]{2: INSERT User (pwd=bcrypt\_hash)};
\draw[ret](8,-2.4)--(4,-2.4)  node[midway,below,font=\tiny]{3: user\_id=7};
\draw[ret](4,-3.0)--(0,-3.0)  node[midway,below,font=\tiny]{4: 201 Created};
% Étape 2 : Login + OTP
\node[note]at(6,-3.7){Phase 2 : Connexion + OTP};
\draw[msg](0,-4.2)--(4,-4.2)  node[midway,above,font=\tiny]{5: POST /auth/login \{email, password\}};
\draw[msg](4,-5.0)--(8,-5.0)  node[midway,above,font=\tiny]{6: SELECT user → vérif. bcrypt};
\draw[msg](4,-5.8)--(12,-5.8) node[midway,above,font=\tiny]{7: send\_otp\_email(code=284719)};
\draw[ret](4,-6.5)--(0,-6.5)  node[midway,below,font=\tiny]{8: 200 \{message:"OTP envoyé"\}};
% Étape 3 : Validation OTP → JWT
\node[note]at(6,-7.2){Phase 3 : Validation + JWT};
\draw[msg](0,-7.8)--(4,-7.8)  node[midway,above,font=\tiny]{9: POST /auth/otp/verify \{email, code\}};
\draw[msg](4,-8.6)--(8,-8.6)  node[midway,above,font=\tiny]{10: SELECT OTP → valide + non-expiré};
\draw[ret](4,-9.3)--(0,-9.3)  node[midway,below,font=\tiny]{11: 200 \{access\_token, refresh\_token\}};
% Étape 4 : Requête authentifiée
\node[note]at(6,-10.0){Phase 4 : Requête authentifiée};
\draw[msg](0,-10.5)--(4,-10.5) node[midway,above,font=\tiny]{12: GET /farms/ (Bearer JWT)};
\draw[msg](4,-11.0)--(8,-11.0) node[midway,above,font=\tiny]{13: SELECT farms WHERE farm\_id=JWT.sub};
\draw[ret](4,-11.5)--(0,-11.5) node[midway,below,font=\tiny]{14: 200 \{farms:[...]\}};
\foreach \x in {0,4,8,12}
  \filldraw[fill=farmblue!15,draw=farmblue!50,rounded corners=1pt]
    (\x-0.15,-0.3) rectangle (\x+0.15,-11.8);
\end{tikzpicture}
\caption{Diagramme de séquence — Authentification JWT + OTP (4 phases)}
\label{fig:seq_auth}
\end{figure}

\section{Sécurité du système}

La sécurité de \farmAI{} v3.0 est conçue selon le principe de \textbf{défense en
profondeur} : chaque couche (authentification, API, IoT, données) dispose de
mécanismes de protection indépendants.

\subsection{Authentification et gestion des identités}

\begin{table}[H]\centering
\caption{Mécanismes d'authentification — Smart Farm AI}
\label{tab:auth_sec}
\begin{tabular}{lll}
\toprule
\textbf{Mécanisme}&\textbf{Implémentation}&\textbf{Protection}\\
\midrule
Hashage mots de passe & bcrypt (coût 12 = $2^{12}$ rounds)  & Brute-force, rainbow tables \\
Token d'accès         & JWT HMAC-SHA256, expiry 30 jours     & Falsification de tokens \\
Token de rafraîchissement & JWT, expiry 7 jours, rotation    & Replay attacks \\
OTP email             & 6 chiffres, 10 min, 3 tentatives max & Credential stuffing \\
\bottomrule
\end{tabular}
\end{table}

\subsection{Sécurité de l'API}

\begin{table}[H]\centering
\caption{Menaces et mitigations — Couche API FastAPI}
\label{tab:api_sec}
\begin{tabular}{llp{5cm}}
\toprule
\textbf{Menace}&\textbf{Mitigation}&\textbf{Implémentation}\\
\midrule
Injection SQL     & SQLAlchemy ORM paramétré  & Requêtes préparées, jamais de \code{text()} brut \\
Injection XSS     & Pydantic validation       & Schémas stricts, HTML auto-escapé \\
CORS non-sécurisé & Liste blanche d'origines  & \code{allow\_origins=["https://farm.ai"]} \\
DDoS / Rate limit & FastAPI slowapi           & 100 req/min par IP, 429 si dépassé \\
HTTPS forcé       & Caddy ACME (prod)         & TLS 1.3, HSTS, certificats auto-renouvelés \\
\bottomrule
\end{tabular}
\end{table}

\subsection{Sécurité IoT}

\begin{tcolorbox}[warningbox,title={Modèle de menace IoT — Surface d'attaque ESP32}]
\begin{description}
  \item[\textbf{Réseau local uniquement}] Les nœuds ESP32 communiquent sur le réseau
        Wi-Fi privé de la ferme. L'endpoint \api{/iot/telemetry} n'est pas exposé à
        Internet (derrière le reverse proxy Caddy).
  \item[\textbf{farm\_id côté serveur}] L'identifiant de la ferme est extrait du JWT
        et jamais fourni par le nœud IoT. Un ESP32 compromis ne peut pas injecter des
        données dans une autre ferme.
  \item[\textbf{Watchdog 30 min}] Redémarrage automatique du firmware si aucun cycle
        \code{loop()} n'est exécuté — protection contre les blocages logiciels.
  \item[\textbf{Limite connue}] HTTP (non HTTPS) entre ESP32 et serveur : acceptable
        en réseau local, vulnérable si exposé à Internet. Migration MQTT+TLS planifiée
        en v4.0.
\end{description}
\end{tcolorbox}

\subsection{Isolation multi-tenant}

\begin{lstlisting}[style=pythonstyle,caption={Isolation multi-tenant — FastAPI Depends},label={lst:tenant}]
# Chaque endpoint injecte automatiquement le farm_id depuis le JWT
async def get_current_farm(
    token: str = Depends(oauth2_scheme),
    db: Session = Depends(get_db)
) -> Farm:
    payload  = decode_jwt(token)
    farm_id  = payload["farm_id"]          # toujours extrait du JWT
    farm     = db.query(Farm).filter(
        Farm.id == farm_id,
        Farm.owner_id == payload["sub"]    # double vérification propriétaire
    ).first()
    if not farm:
        raise HTTPException(403, "Accès refusé")
    return farm
\end{lstlisting}

\noindent\textbf{Conclusion du chapitre :} Ce chapitre a présenté l'architecture
complète de \farmAI{} v3.0 — diagrammes UML (Use Cases, séquences), architecture en
5 couches, Clean Architecture backend, configuration IoT et dispositifs de sécurité.
Le chapitre suivant détaille le pipeline MLOps et les modèles d'IA.

% ============================================================
\chapter{MLOps et Pipeline d'Intelligence Artificielle}
% ============================================================

\noindent Ce chapitre présente l'infrastructure Data Science et MLOps de \farmAI{}
v3.0 en trois axes : l'analyse exploratoire des données (EDA), les fondements
mathématiques des métriques utilisées, et le pipeline industriel DVC+MLflow.
La section finale documente les 8 modèles YOLO, l'IsolationForest et l'assistant
souverain RAG-Darija, avec une justification rigoureuse de chaque choix algorithmique.

\section{Analyse exploratoire des données (EDA)}

\subsection{Dataset IoT — NODE\_A et NODE\_B}

Avant toute modélisation, les données télémétriques ont été analysées pour comprendre
leur distribution, identifier les valeurs manquantes et détecter les anomalies évidentes.

\begin{table}[H]\centering
\caption{Statistiques descriptives des features IoT (NODE\_A + NODE\_B)}
\label{tab:eda_iot}
\small
\begin{tabular}{lccccc}
\toprule
\textbf{Feature}&\textbf{Min}&\textbf{Max}&\textbf{Moyenne}&\textbf{Écart-type}&\textbf{\% manquant}\\
\midrule
soil\_humidity (NODE\_A) & 12.3 & 98.7 & 51.4 & 18.2 & 0.8\% \\
pressure (NODE\_A)       & 0.0  & 12.1 & 4.8  & 2.1  & 1.2\% \\
flow (NODE\_A)           & 0.0  & 8.3  & 2.1  & 1.7  & 3.4\% \\
temperature (NODE\_A)    & 8.2  & 42.1 & 24.7 & 7.3  & 0.3\% \\
weight (NODE\_B)         & 18.0 & 38.5 & 26.3 & 4.1  & 0.5\% \\
hive\_temp (NODE\_B)     & 22.1 & 42.8 & 34.2 & 3.6  & 1.0\% \\
ext\_hum (NODE\_B)       & 28.0 & 95.0 & 62.4 & 14.8 & 0.7\% \\
\midrule
\textbf{Taux anomalies} & \multicolumn{5}{l}{5.2\% détectées par IsolationForest après nettoyage} \\
\bottomrule
\end{tabular}
\end{table}

\noindent\textbf{Observations clés :} La feature \code{flow} présente 3.4\% de valeurs
manquantes (capteur déconnecté lors des nuits sans irrigation). La \code{hive\_temp}
montre une distribution bimodale (jour : $\mu=35.8$°C, nuit : $\mu=32.1$°C) qui
justifie l'usage de l'IsolationForest plutôt qu'un seuil fixe.

\subsection{Datasets YOLO — Composition et équilibre}

\begin{table}[H]\centering
\caption{Composition des datasets d'entraînement YOLO}
\label{tab:eda_yolo}
\small
\begin{tabular}{lccccc}
\toprule
\textbf{Modèle}&\textbf{Train}&\textbf{Val}&\textbf{Test}&\textbf{Classes}&\textbf{Déséquilibre max}\\
\midrule
bee\_obb      & 1\,200 & 240 & 120 & 1   & 1:1 (équilibré) \\
leaf\_disease & 8\,500 & 1\,700 & 850 & 12 & 1:8 (rouille:mildiou) \\
fire\_smoke   & 3\,200 & 640 & 320 & 2   & 1:1.3 \\
livestock     & 2\,800 & 560 & 280 & 4   & 1:3 \\
olive         & 1\,600 & 320 & 160 & 5   & 1:4 \\
insects       & 900    & 180 & 90  & 8   & 1:8 (\textbf{problématique}) \\
lemon         & 2\,100 & 420 & 210 & 4   & 1:2 \\
orange        & 1\,950 & 390 & 195 & 4   & 1:2 \\
\bottomrule
\end{tabular}
\end{table}

\noindent\textbf{Stratégie d'augmentation} pour les classes déséquilibrées (leaf\_disease,
insects) : Mosaic $\times$4 (assembler 4 images en une), RandomFlip horizontal/vertical,
HSV Jitter ($H\pm$0.015, $S\pm$0.7, $V\pm$0.4), CutMix sur les classes sous-représentées.

\subsection{Corpus RAG — Sources documentaires}

\begin{table}[H]\centering
\caption{Corpus de l'assistant souverain (ChromaDB)}
\label{tab:eda_rag}
\begin{tabular}{lccc}
\toprule
\textbf{Source}&\textbf{Documents}&\textbf{Tokens moy.}&\textbf{Langue}\\
\midrule
UTAP — Guides apiculture   & 45 & 850   & AR / FR \\
AVFA — Fiches cultures     & 32 & 1\,200 & FR / Darija \\
Guides régionaux UTAP      & 18 & 600   & Darija \\
\midrule
\textbf{Total}             & \textbf{95} & \textbf{920 moy.} & Trilingue \\
\bottomrule
\end{tabular}
\end{table}

\section{Fondements mathématiques des métriques}

\subsection{Métriques de détection d'objets (YOLO)}

Soient TP les vrais positifs (détections correctes), FP les faux positifs (fausses
détections) et FN les faux négatifs (objets manqués) :

\begin{align}
\text{Précision} &= \frac{TP}{TP + FP} \label{eq:precision}\\[4pt]
\text{Rappel}    &= \frac{TP}{TP + FN}  \label{eq:recall}\\[4pt]
F_1              &= 2 \cdot \frac{\text{Précision} \times \text{Rappel}}{\text{Précision} + \text{Rappel}} \label{eq:f1}
\end{align}

L'\textbf{Average Precision} (AP) par classe $c$ est l'aire sous la courbe
Précision-Rappel :

\begin{equation}
AP_c = \int_0^1 P_c(r)\,dr \;\approx\; \sum_{k} \bigl(R_{k+1} - R_k\bigr) \cdot P_{k+1}
\label{eq:ap}
\end{equation}

Le \textbf{mAP50} est la moyenne des AP sur toutes les $|C|$ classes, calculé avec
un seuil d'Intersection over Union (IoU) de 0.50 :

\begin{equation}
\boxed{
\text{mAP50} = \frac{1}{|C|} \sum_{c \in C} AP_c \;\Big|_{\text{IoU} \geq 0.50}
}
\label{eq:map50}
\end{equation}

Pour \farmAI{} : $|C| = 40$ classes réparties sur 8 modèles. La valeur globale
$\text{mAP50} = 0.891$ signifie que, en moyenne, 89.1\% des détections avec IoU$\geq$0.50
sont correctes.

\subsection{Métriques de détection d'anomalies (IsolationForest)}

Le \textbf{score d'anomalie} de l'IsolationForest pour un point $\mathbf{x}$ sur $n$
observations est :

\begin{equation}
s(\mathbf{x}, n) = 2^{-\frac{\mathbb{E}[h(\mathbf{x})]}{c(n)}}
\label{eq:isof}
\end{equation}

où $h(\mathbf{x})$ est la profondeur d'isolation (nombre de coupures pour isoler $\mathbf{x}$)
et $c(n) = 2H(n-1) - 2(n-1)/n$ est la profondeur moyenne d'un arbre aléatoire
($H$ = nombre harmonique). Si $s \to 1$ : anomalie ; $s \to 0.5$ : point normal.

Le \textbf{Silhouette Score} évalue la qualité de la séparation inliers/anomalies :

\begin{equation}
s(i) = \frac{b(i) - a(i)}{\max\bigl(a(i),\, b(i)\bigr)}, \quad s(i) \in [-1, +1]
\label{eq:silhouette}
\end{equation}

où $a(i) =$ distance intra-cluster moyenne, $b(i) =$ distance au cluster le plus proche.

\begin{equation}
\boxed{
\bar{S} = \frac{1}{n} \sum_{i=1}^n s(i) = 0.73 \quad \Rightarrow \quad \text{Structure forte}
}
\label{eq:sil_result}
\end{equation}

\begin{table}[H]\centering
\caption{Interprétation des métriques — résultats FARM AI}
\label{tab:metrics_interp}
\begin{tabular}{llll}
\toprule
\textbf{Métrique}&\textbf{Formule (réf.)}&\textbf{Valeur FARM AI}&\textbf{Interprétation}\\
\midrule
mAP50  & Éq.~\ref{eq:map50} & 0.891 (global) & 89.1\% de détections correctes \\
mAP50 (bee) & Éq.~\ref{eq:map50} & 0.914 & Excellent \\
Silhouette $\bar{S}$ & Éq.~\ref{eq:silhouette} & 0.73 & Bonne séparation anomalies \\
Anomaly rate & $\frac{\#anomalies}{n}$ & 5.2\% & Cohérent avec contamination=0.05 \\
F1-score moyen & Éq.~\ref{eq:f1} & 0.883 & Précision/Rappel équilibrés \\
\bottomrule
\end{tabular}
\end{table}

\section{Pipeline DVC + MLflow}

\begin{figure}[H]\centering
\begin{tikzpicture}[x=1cm,y=1cm,
  st/.style={rectangle,rounded corners=4pt,draw=#1!70!black,fill=#1!15,
    text width=2.4cm,minimum height=1.3cm,align=center,font=\scriptsize\bfseries,inner sep=3pt},
  art/.style={rectangle,rounded corners=2pt,draw=farmgray!50,fill=farmgray!8,
    text width=2cm,minimum height=0.8cm,align=center,font=\tiny,inner sep=3pt,dashed},
  met/.style={rectangle,rounded corners=2pt,draw=farmgreen!50,fill=farmgreen!8,
    text width=2.2cm,minimum height=0.8cm,align=center,font=\tiny,inner sep=3pt}]

\node[st=farmgray]  (csv)     at(0,0)    {CSV IoT\\NODE\_A/B};
\node[st=farmblue]  (dvc1)    at(3.5,0)  {DVC Stage 1\\build\_dataset};
\node[st=farmblue]  (dvc2)    at(7,0)    {DVC Stage 2\\train\_anomaly};
\node[st=farmgreen] (mlfrun)  at(10.5,0) {MLflow Run\\log metrics};
\node[st=farmgreen] (reg)     at(14,0)   {MLflow\\Registry};
\node[st=farmorange](serve)   at(14,-3)  {serve\_latest\\\_model.py};
\node[st=farmred]   (alert)   at(10.5,-3){Anomaly\\Alert \& WS};

\foreach \a/\b in {csv/dvc1,dvc1/dvc2,dvc2/mlfrun,mlfrun/reg}
  \draw[->,thick,>=Stealth,farmblue!70](\a)--(\b);
\draw[->,thick,>=Stealth,farmgreen!70](reg.south)--(serve.north);
\draw[->,thick,>=Stealth,farmred!70](serve.west)--(alert.east);

\node[art]at(1.75,2){data/raw/\\telemetry.csv};
\node[art]at(5.25,2){data/proc/\\features.csv};
\node[art]at(8.75,2){models/\\anomaly.pkl};
\node[met]at(10.5,-1.7){Sil.=0.73\\Anom.=5\%};
\node[met]at(14,-1.7){TelemetryAnomaly\\Detector v1};

\draw[decorate,decoration={brace,amplitude=5pt},farmblue,thick]
  (0.5,-0.75)--(8.5,-0.75)node[midway,below=7pt,font=\tiny\bfseries\color{farmblue}]
  {Pipeline DVC versionné};
\draw[decorate,decoration={brace,amplitude=5pt},farmgreen,thick]
  (9,-0.75)--(15.5,-0.75)node[midway,below=7pt,font=\tiny\bfseries\color{farmgreen}]
  {MLflow Tracking};
\end{tikzpicture}
\caption{Pipeline MLOps — DVC, MLflow, IsolationForest}
\label{fig:mlops}
\end{figure}

\begin{lstlisting}[style=pythonstyle,caption={train\_anomaly.py — Stage 2 DVC avec MLflow},label={lst:train}]
import mlflow, joblib
from sklearn.ensemble import IsolationForest
from sklearn.metrics import silhouette_score

def main():
    X = pd.read_csv("data/processed/features.csv")[FEATURES].dropna().values

    with mlflow.start_run(run_name="IsolationForest_v3"):
        clf = IsolationForest(n_estimators=150, contamination=0.05,
                              random_state=42, n_jobs=-1)
        clf.fit(X)

        labels = clf.predict(X)          # -1 anomalie, +1 inlier
        sil    = silhouette_score(X, labels)

        mlflow.log_param("n_estimators",  150)
        mlflow.log_param("contamination", 0.05)
        mlflow.log_metric("silhouette_score", round(sil, 4))
        mlflow.log_metric("anomaly_rate", (labels==-1).mean())

        joblib.dump(clf, "models/anomaly_detector.pkl")
        mlflow.sklearn.log_model(clf, "model",
            registered_model_name="TelemetryAnomalyDetector")
\end{lstlisting}

\section{Registre des 8 modèles YOLO v11}

\begin{table}[H]\centering
\caption{Modèles YOLO intégrés dans Smart Farm AI v3.0}
\label{tab:yolo}
\small
\begin{tabular}{llccc}
\toprule
\textbf{Modèle}&\textbf{Domaine / Espèce}&\textbf{Classes}&\textbf{Archi.}&\textbf{mAP50}\\
\midrule
bee\_obb       & Abeilles / ruche OBB    & 1  & YOLOv11-OBB & 0.914\\
livestock      & Bétail (maladies)       & 4  & YOLOv11-det & 0.891\\
leaf\_disease  & Maladies foliaires      & 12 & YOLOv11-det & 0.879\\
olive          & Maladies olivier        & 5  & YOLOv11-det & 0.863\\
insects        & Insectes ravageurs      & 8  & YOLOv11-det & 0.847\\
lemon          & Maladies citronnier     & 4  & YOLOv11-det & 0.901\\
orange         & Maladies oranger        & 4  & YOLOv11-det & 0.895\\
fire\_smoke    & Feu / fumée (sécurité)  & 2  & YOLOv11-det & 0.934\\
\midrule
\textbf{Total} & 8 modèles               &\textbf{40}&YOLOv11&\textbf{0.891}\\
\bottomrule
\end{tabular}
\end{table}

\section{Justification des choix algorithmiques}

La rigueur scientifique exige de justifier chaque choix technologique par rapport aux
alternatives disponibles.

\subsection{Vision par ordinateur : YOLO v11 vs alternatives}

\begin{table}[H]\centering
\caption{Comparaison des architectures de détection d'objets}
\label{tab:yolo_compare}
\small
\begin{tabular}{lccll}
\toprule
\textbf{Modèle}&\textbf{Latence CPU (ms)}&\textbf{mAP COCO}&\textbf{Architecture}&\textbf{Décision}\\
\midrule
\textbf{YOLO v11} & \textbf{167}  & \textbf{53.9} & Single-pass anchor-free & \yes~\textbf{Choisi}\\
YOLOv8            & 185           & 52.9          & Single-pass             & Alternative valide \\
EfficientDet-D2   & 300           & 55.1          & BiFPN + compound scale  & Rejeté (ARM64 lourd)\\
Faster R-CNN      & 500           & 55.7          & Two-stage region prop.  & Rejeté (trop lent)\\
SSD MobileNet     & 120           & 43.1          & Single-pass + anchors   & Rejeté (précision faible)\\
\bottomrule
\end{tabular}
\end{table}

\noindent\textbf{Justification :} YOLO v11 offre le meilleur compromis latence/précision
sur CPU sans GPU. Son architecture \textit{anchor-free} simplifie la configuration des
boîtes englobantes, et sa prise en charge native des \textit{Oriented Bounding Boxes}
(OBB) est indispensable pour la détection des ruches qui apparaissent sous des angles
variables dans les photographies de terrain.

\subsection{Détection d'anomalies : IsolationForest vs alternatives}

\begin{table}[H]\centering
\caption{Comparaison des algorithmes de détection d'anomalies}
\label{tab:anomaly_compare}
\small
\begin{tabular}{lcclll}
\toprule
\textbf{Algorithme}&\textbf{Complexité}&\textbf{Params}&\textbf{Streaming}&\textbf{Interprétabilité}&\textbf{Décision}\\
\midrule
\textbf{IsolationForest} & $O(n \log n)$ & 2   & \yes~Partiel & Score continu & \yes~\textbf{Choisi}\\
LOF                      & $O(n^2)$      & 1 (k) & \no          & Score local    & Rejeté (échelle)\\
Autoencoder              & $O(\text{epochs} \times n)$ & $>$100K & \no & Opaque & Rejeté (complexité)\\
DBSCAN                   & $O(n \log n)$ & 2   & \no~Recluster & Labels discrets & Rejeté (statique)\\
One-Class SVM            & $O(n^2)$      & 2   & \no          & Score kernel   & Rejeté (non-scalable)\\
\bottomrule
\end{tabular}
\end{table}

\noindent\textbf{Justification :} L'IsolationForest est optimal pour les données
télémétriques IoT : complexité sub-quadratique, insensibilité à la malédiction de la
dimensionnalité, et production d'un \textit{score d'anomalie continu} (pas une décision
binaire) permettant de prioriser les alertes par sévérité. Les 2 hyperparamètres
(\code{n\_estimators}, \code{contamination}) sont interprétables métier.

\subsection{LLM Darija : Labess-7B vs alternatives}

\begin{table}[H]\centering
\caption{Comparaison des modèles de langage pour le contexte tunisien}
\label{tab:llm_compare}
\small
\begin{tabular}{lllcll}
\toprule
\textbf{Modèle}&\textbf{Darija natif}&\textbf{Taille}&\textbf{Ollama}&\textbf{Licence}&\textbf{Décision}\\
\midrule
\textbf{Labess-7B} & \yes~Tunisien natif & 7B & \yes & MIT    & \yes~\textbf{Choisi}\\
AraBERT            & \no~MSA seulement  & 135M encodeur & \no & Apache 2.0 & Rejeté (Darija absent)\\
CAMeL-LM           & Partiel (dialectes) & 310M & \no & CC-NC  & Rejeté (non-commercial)\\
AraGPT2            & Partiel             & 1.5B & \no & MIT    & Rejeté (pas Ollama)\\
Llama3.1-8B        & \no~Anglais/FR     & 8B  & \yes & Llama  & Rejeté (Darija absent)\\
\bottomrule
\end{tabular}
\end{table}

\noindent\textbf{Justification :} Labess-7B est le seul LLM natif pour le dialecte
tunisien opérationnel sur Ollama. Son entraînement sur des données conversationnelles
tunisiennes garantit une compréhension des tournures dialectales agricoles que ne
peuvent pas reproduire les LLMs génériques.

\subsection{Versionnement MLOps : DVC + MLflow (complémentaires)}

DVC et MLflow répondent à deux besoins distincts et complémentaires :
\begin{itemize}
  \item \textbf{DVC} = versionnement \textit{données} et \textit{pipelines} (Git pour
        les gros fichiers CSV, .pkl, .pt). Garantit la reproductibilité : la même
        commande \code{dvc repro} produit toujours le même modèle.
  \item \textbf{MLflow} = tracking des \textit{expériences} (hyperparamètres, métriques,
        artefacts). Permet de comparer 10 runs avec différents \code{contamination} en
        un coup d'œil.
\end{itemize}
Les alternatives (W\&B, Neptune, Comet) nécessitent une connexion cloud. DVC+MLflow
fonctionnent entièrement en local, conforme à l'objectif de souveraineté technologique.

\section{Assistant Souverain — RAG Darija}

\begin{tcolorbox}[technote,title={Architecture de l'assistant}]
\begin{enumerate}
  \item \textbf{Labess-7B} (wghezaiel/labess-7b-chat) via Ollama local — LLM Darija tunisien.
  \item \textbf{RAG ChromaDB} — collection \texttt{sovereign\_wisdom\_v3}, HNSW cosine,
        sources UTAP + AVFA encodées avec \texttt{paraphrase-multilingual-mpnet}.
  \item \textbf{LLaVA} — analyse multimodale d'images ESP32-CAM.
  \item \textbf{Fallback} : Ollama local $\to$ Groq API $\to$ base statique (Lite Mode Oracle).
\end{enumerate}
\end{tcolorbox}

\begin{lstlisting}[style=pythonstyle,caption={Service RAG ChromaDB (extrait)},label={lst:rag}]
collection = chroma_client.get_or_create_collection(
    name="sovereign_wisdom_v3",
    metadata={"hnsw:space": "cosine"}
)

def query_rag(question: str, n: int = 5) -> list[str]:
    emb = encoder.encode([question]).tolist()
    return collection.query(
        query_embeddings=emb, n_results=n,
        include=["documents"]
    )["documents"][0]

async def chat(question: str, profile: str) -> str:
    if settings.LITE_MODE:
        return STATIC_KB[profile]["response"]
    context  = query_rag(question)
    prompt   = build_prompt(profile, question, context)
    return ollama_generate(prompt)  # Labess-7B local
\end{lstlisting}

% ============================================================
\chapter{Réalisation et Tests}
% ============================================================

\noindent Ce chapitre documente la concrétisation technique du projet \farmAI{} v3.0.
Il présente l'environnement de développement, les 20 endpoints API principaux, les
interfaces utilisateur clés, puis clôt par une discussion critique des résultats
— incluant l'analyse des limites des modèles YOLO, de l'IsolationForest et du LLM
Darija — fournissant ainsi une évaluation honnête et rigoureuse des performances.

\section{Stack technologique}

\begin{table}[H]\centering
\caption{Environnement de développement}
\label{tab:stack}
\begin{tabular}{lll}
\toprule
\textbf{Composant}&\textbf{Technologie}&\textbf{Version}\\
\midrule
Backend API       & FastAPI + Python    & 0.100+ / 3.11\\
ORM / BDD         & SQLAlchemy + PostgreSQL & 2.0 / 15\\
Vision IA         & Ultralytics YOLO    & 8.1\\
LLM local         & Ollama + Labess-7B  & 0.1.x / 7B\\
RAG               & ChromaDB            & 0.4+\\
Anomalie IoT      & scikit-learn        & 1.4\\
MLOps             & DVC + MLflow        & 3.x / 2.x\\
Frontend          & React + Vite        & 18.2 / 5.1\\
Graphiques        & Recharts            & 2.x\\
Cartographie      & Leaflet             & 1.9\\
PWA offline       & Dexie.js            & 3.x\\
IoT firmware      & PlatformIO (ESP32)  & 6.x\\
Simulation        & Wokwi               & --\\
Déploiement       & Docker Compose      & 2.x\\
Cloud             & Oracle Cloud ARM64  & Always Free\\
\bottomrule
\end{tabular}
\end{table}

\section{Endpoints API principaux}

\begin{longtable}{lll cc}
\caption{Tableau des endpoints FastAPI} \label{tab:endpoints}\\
\toprule
\textbf{Méth.}&\textbf{URL}&\textbf{Description}&\textbf{Auth}&\textbf{Spr.}\\
\midrule\endfirsthead
\toprule\textbf{Méth.}&\textbf{URL}&\textbf{Description}&\textbf{Auth}&\textbf{Spr.}\\\midrule\endhead
\bottomrule\endfoot
POST&\api{/auth/register}       &Création compte + bcrypt         &Public&S1\\
POST&\api{/auth/login}          &JWT (access + refresh)           &Public&S1\\
POST&\api{/auth/otp/send}       &OTP email 6 chiffres             &Public&S1\\
POST&\api{/auth/otp/verify}     &Validation OTP 10 min            &Public&S1\\
GET &\api{/farms/}              &Liste fermes du owner            &JWT  &S1\\
POST&\api{/farms/}              &Créer ferme + GPS                &JWT  &S1\\
POST&\api{/iot/telemetry}       &Réception métrique ESP32         &Public&S2\\
GET &\api{/iot/latest}          &Dernières métriques NODE\_A/B    &JWT  &S2\\
GET &\api{/iot/history}         &Historique avec filtre           &JWT  &S2\\
WS  &\api{/iot/ws/alerts}       &Alertes WebSocket temps réel     &JWT  &S2\\
GET &\api{/bee/yards}           &Liste apiaires                   &JWT  &S3\\
POST&\api{/bee/hives}           &Créer ruche + QR Code            &JWT  &S3\\
POST&\api{/bee/visits}          &Enregistrer visite               &JWT  &S3\\
POST&\api{/bee/harvests}        &Récolte miel                     &JWT  &S3\\
GET &\api{/bee/analytics}       &Stats apicoles                   &JWT  &S3\\
POST&\api{/cv/scan}             &Inférence YOLO (upload)          &JWT  &S4\\
GET &\api{/cv/models}           &8 modèles disponibles            &JWT  &S4\\
GET &\api{/cv/history}          &Historique scans                 &JWT  &S4\\
POST&\api{/assistant/chat}      &Chat Labess-7B + RAG             &JWT  &S5\\
GET &\api{/mlops/metrics}       &Métriques MLflow                 &JWT  &S6\\
\bottomrule
\end{longtable}

\section{Tests et résultats}

\begin{table}[H]\centering
\caption{Résultats des tests backend (pytest)}
\label{tab:tests}
\begin{tabular}{llcc}
\toprule
\textbf{Suite}&\textbf{Cas de test}&\textbf{Résultat}&\textbf{Couverture}\\
\midrule
Authentification  & Register, Login, OTP, refresh  & 12/12 & 87\%\\
Fermes            & CRUD + isolation multi-tenant  & 8/8   & 92\%\\
IoT Télémétrie    & POST, GET latest, anomalie     & 6/6   & 85\%\\
Apiculture        & CRUD ruches, visites, récoltes & 10/10 & 89\%\\
CV Scan YOLO      & Inférence 3 modèles (mock)     & 6/6   & 78\%\\
\midrule
\textbf{Total}    &                                &\textbf{42/42}&\textbf{86\%}\\
\bottomrule
\end{tabular}
\end{table}

\begin{tcolorbox}[infobox,title={Performance YOLO sur CPU Intel Core i7}]
\begin{itemize}
  \item Latence inférence : \textbf{167 ms} (p50) / \textbf{198 ms} (p95) — SLA $<$200 ms \yes
  \item Modèle le plus lent : leaf\_disease (192 ms — 12 classes)
  \item Modèle le plus rapide : fire\_smoke (124 ms — 2 classes)
  \item Simulation Wokwi : 100\% des métriques reçues côté backend
\end{itemize}
\end{tcolorbox}

\section{Discussion critique des résultats}

\subsection{Analyse des résultats YOLO}

Les huit modèles YOLO v11 présentent des performances hétérogènes qui méritent une
analyse critique :

\begin{table}[H]\centering
\caption{Analyse critique des résultats YOLO par domaine}
\label{tab:yolo_discuss}
\small
\begin{tabular}{llcp{5.5cm}}
\toprule
\textbf{Modèle}&\textbf{mAP50}&\textbf{Verdict}&\textbf{Analyse}\\
\midrule
fire\_smoke    & 0.934 & \yes~Excellent  & Classes visuellement distinctives (rouge/gris), dataset équilibré \\
bee\_obb       & 0.914 & \yes~Excellent  & OBB efficace sur ruches rectangulaires orientées \\
lemon / orange & 0.901/0.895 & \yes~Bon & Symétrie visuelle des maladies agrumicoles \\
livestock      & 0.891 & \yes~Bon        & 4 pathologies distinctes visuellement \\
leaf\_disease  & 0.879 & $\sim$ Acceptable & 12 classes créent confusions entre maladies similaires \\
olive          & 0.863 & $\sim$ Acceptable & Dataset limité pour pathologies locales tunisiennes \\
insects        & 0.847 & \textcolor{farmorange}{Moyen} & \textbf{Faiblesse principale} : espèces visuellement similaires, dataset déséquilibré ×8 \\
\bottomrule
\end{tabular}
\end{table}

\begin{tcolorbox}[warningbox,title={Limite YOLO : modèle insects (mAP50 = 0.847)}]
Le modèle de détection d'insectes ravageurs présente la performance la plus faible.
Cause principale : les datasets publics disponibles sont principalement composés d'images
européennes et américaines, peu représentatives des espèces ravageuses tunisiennes
(thrips, cochenilles, pucerons locaux). \textbf{Recommandation :} constituer un dataset
photographique local avec la collaboration de l'AVFA (Agence de Vulgarisation et de
Formation Agricoles).
\end{tcolorbox}

\subsection{Analyse du modèle IsolationForest}

Le score de Silhouette obtenu ($\bar{S} = 0.73$) est qualifié de \textbf{bon} selon
les seuils standards (Kaufman \& Rousseeuw, 1990) :

\begin{center}
\begin{tabular}{cll}
\toprule
\textbf{Plage Silhouette}&\textbf{Interprétation}&\textbf{FARM AI}\\
\midrule
$[0.71, 1.00]$ & Structure forte       & \yes~$\bar{S} = 0.73$ \\
$[0.51, 0.70]$ & Structure raisonnable & -- \\
$[0.26, 0.50]$ & Structure faible      & -- \\
$[-1,  0.25]$  & Pas de structure      & -- \\
\bottomrule
\end{tabular}
\end{center}

\begin{tcolorbox}[warningbox,title={Limites de l'IsolationForest sur données IoT}]
\begin{itemize}
  \item \textbf{Distribution bimodale jour/nuit} : les lectures d'humidité sol et de
        débit diffèrent systématiquement entre irrigations nocturnes et diurnes. Cela
        crée une zone ambiguë à la frontière inlier/anomalie (scores proches de 0.5).
  \item \textbf{Modèle statique} : pas de réentraînement incrémental. Un changement de
        capteur ou une nouvelle saison peut dériver la distribution sans adaptation.
  \item \textbf{Calibration du paramètre contamination} : la valeur fixée à 0.05 est
        une hypothèse \textit{a priori}. Une validation croisée temporelle serait
        nécessaire pour confirmer ce taux sur données réelles de terrain.
\end{itemize}
\textbf{Recommandation :} intégrer un pipeline de réentraînement mensuel automatique
via DVC + GitHub Actions pour adapter le modèle aux nouvelles saisons.
\end{tcolorbox}

\subsection{Limites du LLM Darija (Labess-7B)}

\begin{tcolorbox}[warningbox,title={Limites de l'assistant souverain}]
\begin{itemize}
  \item \textbf{Latence CPU (15--25 s)} : incompatible avec une interaction mobile
        temps réel. Le fallback Groq API réduit ce délai à $<$3 s mais introduit une
        dépendance cloud externe contraire à l'objectif de souveraineté.
  \item \textbf{Corpus RAG limité (95 documents)} : couvre les pathologies majeures
        mais insuffisant pour les cas rares (maladies émergentes, nouvelles espèces
        invasives). Aucune métrique BLEU ou ROUGE n'a été formellement calculée sur la
        qualité des réponses.
  \item \textbf{Dialecte Darija non uniforme} : la Darija varie significativement entre
        régions tunisiennes (Tunis, Sfax, Sousse). Labess-7B est biaisé vers le dialecte
        de la capitale.
\end{itemize}
\end{tcolorbox}

\subsection{Limites de la simulation IoT}

\begin{tcolorbox}[warningbox,title={Wokwi vs matériel réel}]
\begin{itemize}
  \item \textbf{Capteurs virtuels déterministes} : Wokwi simule des capteurs parfaits
        sans bruit électronique, dérive thermique ni erreurs de lecture. Les capteurs
        réels présentent $\pm$2\% d'imprécision systématique nécessitant un filtre
        passe-bas (ex. filtre de Kalman).
  \item \textbf{Absence de chiffrement HTTP} : le protocole HTTP entre ESP32 et
        serveur est acceptable en réseau local privé, mais constitue une vulnérabilité
        si le déploiement est exposé à Internet. La migration vers HTTPS (MQTT+TLS)
        est requise pour un déploiement public.
  \item \textbf{Fréquence 10 s insuffisante} pour la détection de vibrations des ruches
        (fréquences 200--300 Hz caractéristiques de l'essaimage, nécessitent $\geq$600 Hz
        d'échantillonnage selon le théorème de Nyquist).
\end{itemize}
\end{tcolorbox}

\medskip
\noindent\textbf{Conclusion du chapitre :} Ce chapitre a présenté l'environnement de
développement, les interfaces clés et les résultats de tests. L'analyse critique révèle
des performances globalement satisfaisantes (mAP50 moyen 0.891, Silhouette 0.73,
42/42 tests) avec des axes d'amélioration identifiés : dataset insects, réentraînement
IsolationForest, et passage HTTP→HTTPS pour le déploiement IoT public.

% ============================================================
\chapter{Conclusion et Perspectives}
% ============================================================

\section{Synthèse des réalisations}

\begin{table}[H]\centering
\caption{Livrables réalisés par sprint}
\label{tab:livr}
\begin{tabular}{clll}
\toprule
\textbf{Spr.}&\textbf{Titre}&\textbf{Livrables clés}&\textbf{Statut}\\
\midrule
S1&Infra \& Auth      &FastAPI 50+ endpoints, JWT+OTP, PostgreSQL 30+ tables&\yes\\
S2&IoT \& Télémétrie  &NODE\_A+B firmware, Wokwi, dashboard Recharts       &\yes\\
S3&Module Apiculture  &HiveIntel 7 modules, CRUD, QR Codes, télémétrie       &\yes\\
S4&Vision IA          &8 YOLO v11, Scanner PWA 13 modèles, Map SIG           &\yes\\
S5&Assistant Souverain&Labess-7B RAG Darija, STT/TTS, Lite Mode Oracle       &\yes\\
S6&MLOps \& Final     &DVC+MLflow (Sil.=0.73), Docker, Oracle Cloud          &\yes\\
\bottomrule
\end{tabular}
\end{table}

\section{Contributions originales}

\begin{description}
  \item[\textbf{Souveraineté IA en Darija :}] Premier assistant conversationnel
        agricole en dialecte tunisien, opérant entièrement en local sans API cloud.
  \item[\textbf{Chaîne IoT-IA complète :}] Même firmware pour simulation (Wokwi)
        et production réelle, avec détection d'anomalies automatique en temps réel.
  \item[\textbf{MLOps agricole reproductible :}] Application des meilleures pratiques
        industrielles (DVC + MLflow) à un contexte agricole tunisien, sans GPU coûteux.
\end{description}

\section{Difficultés rencontrées}

\begin{enumerate}
  \item Déséquilibre des datasets YOLO pour pathologies locales — résolu par augmentation aggressive.
  \item Latence LLM sur CPU (15-25 s) — mitigation via Groq API fallback et Lite Mode.
  \item Oracle Cloud ARM64 sans GPU — LITE\_MODE désactive YOLO et LLM, active KB statique.
  \item Multi-tenancy stricte — propagation systématique du \code{farm\_id} dans tous les modèles SQLAlchemy.
\end{enumerate}

\section{Perspectives d'évolution}

\begin{tcolorbox}[infobox,title={Feuille de route v4.0}]
\textbf{Court terme (3-6 mois) :}
\begin{itemize}
  \item Microphones ruches — analyse fréquences 200-300 Hz (détection essaimage)
  \item Migration TimescaleDB pour séries temporelles IoT
  \item React Native pour Worker (Android/iOS natif)
\end{itemize}
\textbf{Moyen terme (6-18 mois) :}
\begin{itemize}
  \item ESP32-CAM — photos automatiques ruches toutes les heures
  \item Fine-tuning Labess-7B sur corpus agricole tunisien spécialisé
  \item MLOps niveau 2 : CI/CD ML avec GitHub Actions
\end{itemize}
\end{tcolorbox}

\section{Conclusion générale}

\farmAI{} v3.0 démontre qu'il est possible de construire une plateforme IA agricole
complète, souveraine et économiquement accessible en combinant des technologies
open-source de pointe (YOLO v11, FastAPI, ChromaDB, DVC, MLflow) avec une architecture
rigoureuse et des pratiques MLOps industrielles.

La spécificité tunisienne — assistant en Darija, corpus RAG ancré dans les références
officielles nationales, IoT adapté aux conditions locales — fait de \farmAI{} une
fondation réelle pour la transformation numérique de l'agriculture tunisienne.

\begin{tcolorbox}[technote,title={Impact attendu}]
Si 1\% des 516\,000 exploitants tunisiens adoptent \farmAI{} et réduisent leurs pertes
de 15\% grâce à la détection précoce des maladies, cela représente
\textbf{774 fermes bénéficiaires} et une économie estimée à \textbf{2,3 millions TND/an}.
\end{tcolorbox}

% ============================================================
%  BIBLIOGRAPHIE (intégrée — pas besoin de BibTeX)
% ============================================================
\begin{thebibliography}{99}

\bibitem{redmon2016yolo}
J. Redmon, S. Divvala, R. Girshick, A. Farhadi.
\textit{You Only Look Once: Unified, Real-Time Object Detection.}
CVPR, pp. 779--788, 2016.

\bibitem{ultralytics2023yolov8}
G. Jocher, A. Chaurasia, J. Qiu.
\textit{Ultralytics YOLOv8: A State-of-the-Art Object Detection Framework.}
\url{https://github.com/ultralytics/ultralytics}, 2023.

\bibitem{liu2008isolation}
F. T. Liu, K. M. Ting, Z.-H. Zhou.
\textit{Isolation Forest.}
IEEE ICDM, pp. 413--422, 2008.

\bibitem{lecun2015deep}
Y. LeCun, Y. Bengio, G. Hinton.
\textit{Deep Learning.}
Nature, 521(7553), 436--444, 2015.

\bibitem{lewis2020rag}
P. Lewis et al.
\textit{Retrieval-Augmented Generation for Knowledge-Intensive NLP Tasks.}
NeurIPS, vol. 33, pp. 9459--9474, 2020.

\bibitem{labess7b}
W. Ghezaiel.
\textit{Labess-7B: LLM for Tunisian Arabic (Darija).}
\url{https://huggingface.co/wghezaiel/labess-7b-chat}, 2024.

\bibitem{chromadb2023}
\textit{Chroma: The AI-Native Open-Source Embedding Database.}
\url{https://www.trychroma.com}, 2023.

\bibitem{ollama2023}
\textit{Ollama: Run Large Language Models Locally.}
\url{https://ollama.com}, 2023.

\bibitem{fastapi2023}
S. Ramírez.
\textit{FastAPI: Modern, Fast Web Framework for Building APIs with Python.}
\url{https://fastapi.tiangolo.com}, 2023.

\bibitem{mlflow2023}
M. Zaharia et al.
\textit{MLflow: An Open Source Platform for the ML Lifecycle.}
\url{https://mlflow.org}, 2023.

\bibitem{dvc2023}
\textit{DVC: Data Version Control — Git for Data \& Models.}
\url{https://dvc.org}, 2023.

\bibitem{sculley2015hidden}
D. Sculley et al.
\textit{Hidden Technical Debt in Machine Learning Systems.}
NeurIPS, vol. 28, 2015.

\bibitem{crisp_dm}
R. Wirth, J. Hipp.
\textit{CRISP-DM: Towards a Standard Process Model for Data Mining.}
PKDD, pp. 29--39, 2000.

\bibitem{schwaber2020scrum}
K. Schwaber, J. Sutherland.
\textit{The Scrum Guide.}
\url{https://scrumguides.org}, 2020.

\bibitem{beck2001agile}
K. Beck et al.
\textit{Manifesto for Agile Software Development.}
\url{https://agilemanifesto.org}, 2001.

\bibitem{kamilaris2018deep}
A. Kamilaris, F. X. Prenafeta-Boldú.
\textit{Deep Learning in Agriculture: A Survey.}
Computers and Electronics in Agriculture, 147, 70--90, 2018.

\bibitem{liakos2018machine}
K. G. Liakos et al.
\textit{Machine Learning in Agriculture: A Review.}
Sensors, 18(8), 2674, 2018.

\bibitem{ins_tunisie_2023}
Institut National de la Statistique (INS).
\textit{Annuaire Statistique de la Tunisie 2022--2023.}
\url{https://www.ins.tn}, 2023.

\bibitem{fao_tunisie_2022}
FAO.
\textit{Tunisie: Note de Cadrage Agriculture.}
\url{https://www.fao.org/tunisia/fr/}, 2022.

\bibitem{utap2023}
UTAP.
\textit{Guide Technique de l'Apiculture en Tunisie.}
Union Tunisienne de l'Agriculture et de la Pêche, 2023.

\bibitem{espressif2023esp32}
Espressif Systems.
\textit{ESP32: Feature-Rich MCU with Wi-Fi and Bluetooth.}
\url{https://www.espressif.com/en/products/socs/esp32}, 2023.

\bibitem{docker2023}
\textit{Docker: Empowering App Development for Developers.}
\url{https://www.docker.com}, 2023.

\end{thebibliography}

\end{document}
